\documentclass{amsart}
% For dealing with references we use the comment environment
\usepackage{verbatim}
\newenvironment{reference}{\comment}{\endcomment}
%\newenvironment{reference}{}{}
\newenvironment{slogan}{\comment}{\endcomment}
\newenvironment{history}{\comment}{\endcomment}

% For commutative diagrams we use Xy-pic
\usepackage[all]{xy}

% We use 2cell for 2-commutative diagrams.
\xyoption{2cell}
\UseAllTwocells

% We use multicol for the list of chapters between chapters
\usepackage{multicol}

% This is generally recommended for better output
\usepackage{lmodern}
\usepackage[T1]{fontenc}

% For cross-file-references

% Package for hypertext links:
\usepackage{hyperref}

% For any local file, say "hello.tex" you want to link to please

% Theorem environments.
%
\theoremstyle{plain}
\newtheorem{theorem}[subsection]{Theorem}
\newtheorem{proposition}[subsection]{Proposition}
\newtheorem{lemma}[subsection]{Lemma}

\theoremstyle{definition}
\newtheorem{definition}[subsection]{Definition}
\newtheorem{example}[subsection]{Example}
\newtheorem{exercise}[subsection]{Exercise}
\newtheorem{situation}[subsection]{Situation}

\theoremstyle{remark}
\newtheorem{remark}[subsection]{Remark}
\newtheorem{remarks}[subsection]{Remarks}

\numberwithin{equation}{subsection}

% Macros
%
\def\lim{\mathop{\mathrm{lim}}\nolimits}
\def\colim{\mathop{\mathrm{colim}}\nolimits}
\def\Spec{\mathop{\mathrm{Spec}}}
\def\Hom{\mathop{\mathrm{Hom}}\nolimits}
\def\Ext{\mathop{\mathrm{Ext}}\nolimits}
\def\SheafHom{\mathop{\mathcal{H}\!\mathit{om}}\nolimits}
\def\SheafExt{\mathop{\mathcal{E}\!\mathit{xt}}\nolimits}
\def\Sch{\mathit{Sch}}
\def\Mor{\mathop{\mathrm{Mor}}\nolimits}
\def\Ob{\mathop{\mathrm{Ob}}\nolimits}
\def\Sh{\mathop{\mathit{Sh}}\nolimits}
\def\NL{\mathop{N\!L}\nolimits}
\def\CH{\mathop{\mathrm{CH}}\nolimits}
\def\proetale{{pro\text{-}\acute{e}tale}}
\def\etale{{\acute{e}tale}}
\def\QCoh{\mathit{QCoh}}
\def\Ker{\mathop{\mathrm{Ker}}}
\def\Im{\mathop{\mathrm{Im}}}
\def\Coker{\mathop{\mathrm{Coker}}}
\def\Coim{\mathop{\mathrm{Coim}}}

% Boxtimes
%
\DeclareMathSymbol{\boxtimes}{\mathbin}{AMSa}{"02}

%
% Macros for moduli stacks/spaces
%
\def\QCohstack{\mathcal{QC}\!\mathit{oh}}
\def\Cohstack{\mathcal{C}\!\mathit{oh}}
\def\Spacesstack{\mathcal{S}\!\mathit{paces}}
\def\Quotfunctor{\mathrm{Quot}}
\def\Hilbfunctor{\mathrm{Hilb}}
\def\Curvesstack{\mathcal{C}\!\mathit{urves}}
\def\Polarizedstack{\mathcal{P}\!\mathit{olarized}}
\def\Complexesstack{\mathcal{C}\!\mathit{omplexes}}
% \Pic is the operator that assigns to X its picard group, usage \Pic(X)
% \Picardstack_{X/B} denotes the Picard stack of X over B
% \Picardfunctor_{X/B} denotes the Picard functor of X over B
\def\Pic{\mathop{\mathrm{Pic}}\nolimits}
\def\Picardstack{\mathcal{P}\!\mathit{ic}}
\def\Picardfunctor{\mathrm{Pic}}
\def\Deformationcategory{\mathcal{D}\!\mathit{ef}}

\setlength{\emergencystretch}{3em}
\begin{document}
\title[Stacks Project, Tag 08IF]{386 --- Relative Ext with proper flat coherent coefficients is represented by pseudo-coherent complexes}
\maketitle
\noindent Copyright (C) 2005--2025 Johan de Jong. Stacks Project authors. Source: \href{https://stacks.math.columbia.edu/tag/08IF}{Tag 08IF}.

\noindent Editorial difficulty note (not author text). Difficulty H3: The proof must represent all coefficient-variable Ext functors, not only each fixed Ext group. Perfect approximations give finite-degree representing complexes; their duals must then be assembled into a homotopy colimit with controlled truncations to obtain one pseudo-coherent object representing every degree and its boundary maps..

\noindent Editorial provenance note (not author text). History: excerpt from revision \texttt{a0444\allowbreak 6e57e\allowbreak c1fbc\allowbreak 252a8\allowbreak 71afc\allowbreak ec775\allowbreak 2fb28\allowbreak 07b14}, perfect.tex, lines 7057--7180. Reference presentation was adapted; original-author context and core support blocks were retained or added. The mathematical wording is unchanged. Licensed under GNU FDL 1.2 or later; no invariant sections or cover texts. See ../licenses/COPYING. Context blocks reproduce author definitions and supporting results; they are not additional counted problems.

\begin{definition}
\label{algebra-definition-module-finite-type}
Let $R$ be a ring. Let $M$ be an $R$-module.
\begin{enumerate}
\item We say $M$ is a {\it finite $R$-module}, or a {\it finitely generated
$R$-module} if there exist $n \in \mathbf{N}$ and $x_1, \ldots, x_n \in M$
such that every element of $M$ is an $R$-linear combination of the $x_i$.
Equivalently, this means there exists a surjection
$R^{\oplus n} \to M$ for some $n \in \mathbf{N}$.
\item We say $M$ is a {\it finitely presented $R$-module} or an
{\it $R$-module of finite presentation} if there exist integers
$n, m \in \mathbf{N}$ and an exact sequence
$$
R^{\oplus m} \longrightarrow R^{\oplus n} \longrightarrow M \longrightarrow 0
$$
\end{enumerate}
\end{definition}

\begin{definition}
\label{algebra-definition-localization-module}
The $S^{-1}A$-module $S^{-1}M$ is called the {\it localization} of $M$
with respect to $S$.
\end{definition}

\begin{definition}
\label{cohomology-definition-perfect}
Let $(X, \mathcal{O}_X)$ be a ringed space.
Let $\mathcal{E}^\bullet$ be a complex of $\mathcal{O}_X$-modules.
We say $\mathcal{E}^\bullet$ is {\it perfect} if there exists
an open covering $X = \bigcup U_i$ such that for each $i$
there exists a morphism of complexes
$\mathcal{E}_i^\bullet \to \mathcal{E}^\bullet|_{U_i}$
which is a quasi-isomorphism with $\mathcal{E}_i^\bullet$
a strictly perfect complex of $\mathcal{O}_{U_i}$-modules.
An object $E$ of $D(\mathcal{O}_X)$ is {\it perfect}
if it can be represented by a perfect complex of $\mathcal{O}_X$-modules.
\end{definition}

\begin{definition}
\label{cohomology-definition-pseudo-coherent}
Let $(X, \mathcal{O}_X)$ be a ringed space. Let $\mathcal{E}^\bullet$
be a complex of $\mathcal{O}_X$-modules. Let $m \in \mathbf{Z}$.
\begin{enumerate}
\item We say $\mathcal{E}^\bullet$ is {\it $m$-pseudo-coherent}
if there exists an open covering $X = \bigcup U_i$ and for each $i$
a morphism of complexes
$\alpha_i : \mathcal{E}_i^\bullet \to \mathcal{E}^\bullet|_{U_i}$
where $\mathcal{E}_i^\bullet$ is strictly perfect on $U_i$ and
$H^j(\alpha_i)$ is an isomorphism for $j > m$ and $H^m(\alpha_i)$
is surjective.
\item We say $\mathcal{E}^\bullet$ is {\it pseudo-coherent}
if it is $m$-pseudo-coherent for all $m$.
\item We say an object $E$ of $D(\mathcal{O}_X)$ is
{\it $m$-pseudo-coherent} (resp.\ {\it pseudo-coherent})
if and only if it can be represented by a $m$-pseudo-coherent
(resp.\ pseudo-coherent) complex of $\mathcal{O}_X$-modules.
\end{enumerate}
\end{definition}

\begin{definition}
\label{derived-definition-cone}
Let $\mathcal{A}$ be an additive category.
Let $f : K^\bullet \to L^\bullet$ be a morphism of
complexes of $\mathcal{A}$. The {\it cone} of $f$
is the complex $C(f)^\bullet$ given by
$C(f)^n = L^n \oplus K^{n + 1}$ and
differential
$$
d_{C(f)}^n =
\left(
\begin{matrix}
d^n_L & f^{n + 1} \\
0 & -d_K^{n + 1}
\end{matrix}
\right)
$$
It comes equipped with canonical morphisms of complexes
$i : L^\bullet \to C(f)^\bullet$ and $p : C(f)^\bullet \to K^\bullet[1]$
induced by the obvious maps $L^n \to C(f)^n \to K^{n + 1}$.
\end{definition}

\begin{definition}
\label{definition-approximation-holds}
Let $X$ be a scheme. Consider triples $(T, E, m)$ where
\begin{enumerate}
\item $T \subset X$ is a closed subset,
\item $E$ is an object of $D_\QCoh(\mathcal{O}_X)$, and
\item $m \in \mathbf{Z}$.
\end{enumerate}
We say {\it approximation holds for the triple} $(T, E, m)$ if
there exists a perfect object $P$ of $D(\mathcal{O}_X)$ supported on $T$
and a map $\alpha : P \to E$ which induces isomorphisms $H^i(P) \to H^i(E)$
for $i > m$ and a surjection $H^m(P) \to H^m(E)$.
\end{definition}

\begin{definition}
\label{definition-approximation}
Let $X$ be a scheme. We say {\it approximation by perfect complexes holds}
on $X$ if for any closed subset $T \subset X$ with $X \setminus T$
retro-compact in $X$ there exists an integer $r$ such that
for every triple $(T, E, m)$ as in
Definition \ref{definition-approximation-holds} with
\begin{enumerate}
\item $E$ is $(m - r)$-pseudo-coherent, and
\item $H^i(E)$ is supported on $T$ for $i \geq m - r$
\end{enumerate}
approximation holds.
\end{definition}

\begin{situation}
\label{situation-complex}
Here $A$ is a ring and $f_1, \ldots, f_r$ is a sequence of elements of $A$.
We set $X = \Spec(A)$ and $U = D(f_1) \cup \ldots \cup D(f_r) \subset X$.
We denote $\mathcal{U} : U = \bigcup_{i = 1, \ldots, r} D(f_i)$ the
given open covering of $U$.
\end{situation}

\begin{definition}
\label{more-algebra-definition-koszul-complex}
Let $R$ be a ring and let $f_1, \ldots, f_r \in R$. The
{\it Koszul complex on $f_1, \ldots, f_r$} is the Koszul complex
associated to the map $(f_1, \ldots, f_r) : R^{\oplus r} \to R$.
Notation $K_\bullet(f_\bullet)$, $K_\bullet(f_1, \ldots, f_r)$,
$K_\bullet(R, f_1, \ldots, f_r)$, or $K_\bullet(R, f_\bullet)$.
\end{definition}

\begin{lemma}
\label{lemma-lift-pseudo-coherent}
Let $X$ be an affine scheme and let $U \subset X$ be a quasi-compact
open subscheme. For any pseudo-coherent object $E$ of $D(\mathcal{O}_U)$
there exists a bounded above complex of finite free $\mathcal{O}_X$-modules 
whose restriction to $U$ is isomorphic to $E$.
\end{lemma}

\begin{proof}
By Lemma \ref{lemma-pseudo-coherent} we see that $E$ is an object of
$D_\QCoh(\mathcal{O}_U)$. By
Lemma \ref{lemma-lift-quasi-coherent}
we may assume $E = E'|U$ for some object $E'$ of
$D_\QCoh(\mathcal{O}_X)$.
Write $X = \Spec(A)$. By Lemma \ref{lemma-affine-compare-bounded}
we can find a complex $M^\bullet$ of $A$-modules whose associated
complex of $\mathcal{O}_X$-modules is a representative of $E'$.

\medskip\noindent
Choose $f_1, \ldots, f_r \in A$ such that $U = D(f_1) \cup \ldots \cup D(f_r)$.
By Lemma \ref{lemma-pseudo-coherent-affine} the complexes
$M^\bullet_{f_j}$ are pseudo-coherent complexes of $A_{f_j}$-modules.
Let $n$ be an integer. Assume we have a map of complexes
$\alpha : F^\bullet \to M^\bullet$ where $F^\bullet$ is
bounded above, $F^i = 0$ for $i < n$, each $F^i$ is a finite free
$R$-module, such that
$$
H^i(\alpha_{f_j}) : H^i(F^\bullet_{f_j}) \to H^i(M^\bullet_{f_j})
$$
is an isomorphism for $i > n$ and surjective for $i = n$. Picture
$$
\xymatrix{
& F^n \ar[r] \ar[d]^\alpha & F^{n + 1} \ar[d]^\alpha \ar[r] & \ldots \\
M^{n-1} \ar[r] & M^n \ar[r] & M^{n + 1} \ar[r] & \ldots
}
$$
Since each $M^\bullet_{f_j}$ has vanishing cohomology
in large degrees we can find such a map for $n \gg 0$.
By induction on $n$ we are going to extend this to a map
of complexes $F^\bullet \to M^\bullet$
such that $H^i(\alpha_{f_j})$ is an isomorphism
for all $i$. The lemma will follow by taking $\widetilde{F^\bullet}$.

\medskip\noindent
The induction step will be to extend the diagram
above by adding $F^{n - 1}$. Let $C^\bullet$ be the cone on $\alpha$
(Derived Categories, Definition \ref{derived-definition-cone}).
The long exact sequence of cohomology shows that
$H^i(C^\bullet_{f_j}) = 0$ for $i \geq n$. By
More on Algebra, Lemma \ref{more-algebra-lemma-cone-pseudo-coherent}
we see that $C^\bullet_{f_j}$ is $(n - 1)$-pseudo-coherent. By
More on Algebra, Lemma \ref{more-algebra-lemma-finite-cohomology}
we see that $H^{n - 1}(C^\bullet_{f_j})$ is a finite $A_{f_j}$-module.
Choose a finite free $A$-module $F^{n - 1}$ and an $A$-module
$\beta : F^{n - 1} \to C^{n - 1}$ such that the composition
$F^{n - 1} \to C^{n - 1} \to C^n$ is zero and such that
$F^{n - 1}_{f_j}$ surjects onto $H^{n - 1}(C^\bullet_{f_j})$.
(Some details omitted; hint: clear denominators.)
Since $C^{n - 1} = M^{n - 1} \oplus F^n$
we can write $\beta = (\alpha^{n - 1}, -d^{n - 1})$. The vanishing of the
composition $F^{n - 1} \to C^{n - 1} \to C^n$ implies
these maps fit into a morphism of complexes
$$
\xymatrix{
& F^{n - 1} \ar[d]^{\alpha^{n - 1}} \ar[r]_{d^{n - 1}} &
F^n \ar[r] \ar[d]^\alpha &
F^{n + 1} \ar[d]^\alpha \ar[r] & \ldots \\
\ldots \ar[r] &
M^{n - 1} \ar[r] & M^n \ar[r] & M^{n + 1} \ar[r] & \ldots
}
$$
Moreover, these maps define a morphism of distinguished triangles
$$
\xymatrix{
(F^n \to \ldots) \ar[r] \ar[d] &
(F^{n-1} \to \ldots) \ar[r] \ar[d] &
F^{n-1} \ar[r] \ar[d]_\beta &
(F^n \to \ldots)[1] \ar[d] \\
(F^n \to \ldots) \ar[r] &
M^\bullet \ar[r] &
C^\bullet \ar[r] &
(F^n \to \ldots)[1]
}
$$
Hence our choice of $\beta$ implies that the map of complexes
$(F^{-1} \to \ldots) \to M^\bullet$ induces an isomorphism on
cohomology localized at $f_j$ in degrees $\geq n$ and a surjection
in degree $n - 1$. This finishes the proof of the lemma.
\end{proof}


\begin{lemma}
\label{lemma-ext-from-perfect-into-bounded-QCoh}
Let $X$ be a quasi-compact and quasi-separated scheme.
Let $K$ be a perfect object of $D(\mathcal{O}_X)$. Then
\begin{enumerate}
\item there exist integers $a \leq b$ such that for any
$L \in D_\QCoh(\mathcal{O}_X)$ with $H^i(L) = 0$ for $i \in [a, b]$
we have $\Hom_{D(\mathcal{O}_X)}(K, L) = 0$, and
\item if $L$ is bounded, then $\Ext^n_{D(\mathcal{O}_X)}(K, L)$
is zero for all but finitely many $n$.
\end{enumerate}
\end{lemma}

\begin{proof}
Part (2) follows from (1) as $\Ext^n_{D(\mathcal{O}_X)}(K, L) =
\Hom_{D(\mathcal{O}_X)}(K, L[n])$. We prove (1).
Since $K$ is perfect we have
$$
\Hom_{D(\mathcal{O}_X)}(K, L) =
H^0(X, K^\vee \otimes_{\mathcal{O}_X}^\mathbf{L} L)
$$
where $K^\vee$ is the ``dual'' perfect complex to $K$, see
Cohomology, Lemma \href{https://stacks.math.columbia.edu/tag/08DQ}{lemma dual perfect complex (Tag 08DQ)}.
Note that $K^\vee \otimes_{\mathcal{O}_X}^\mathbf{L} L$
is in $D_\QCoh(X)$ by
Lemmas \href{https://stacks.math.columbia.edu/tag/08DX}{lemma quasi coherence tensor product (Tag 08DX)} and
\ref{lemma-pseudo-coherent} (to see that a perfect complex
has quasi-coherent cohomology sheaves). Say $K^\vee$ has
tor amplitude in $[a, b]$. Then the spectral sequence
$$
E_1^{p, q} = H^p(K^\vee \otimes_{\mathcal{O}_X}^\mathbf{L} H^q(L))
\Rightarrow
H^{p + q}(K^\vee \otimes_{\mathcal{O}_X}^\mathbf{L} L)
$$
shows that $H^j(K^\vee \otimes_{\mathcal{O}_X}^\mathbf{L} L)$
is zero if $H^q(L) = 0$ for $q \in [j - b, j - a]$.
Let $N$ be the integer $d$ of Cohomology of Schemes,
Lemma \href{https://stacks.math.columbia.edu/tag/071L}{lemma vanishing nr affines quasi separated (Tag 071L)}.
Then $H^0(X, K^\vee \otimes_{\mathcal{O}_X}^\mathbf{L} L)$
vanishes if the cohomology sheaves
$$
H^{-N}(K^\vee \otimes_{\mathcal{O}_X}^\mathbf{L} L),
\ H^{-N + 1}(K^\vee \otimes_{\mathcal{O}_X}^\mathbf{L} L),
\ \ldots,
\ H^0(K^\vee \otimes_{\mathcal{O}_X}^\mathbf{L} L)
$$
are zero. Namely, by the lemma cited and
Lemma \ref{lemma-application-nice-K-injective}, we have
$$
H^0(X, K^\vee \otimes_{\mathcal{O}_X}^\mathbf{L} L) =
H^0(X, \tau_{\geq -N}(K^\vee \otimes_{\mathcal{O}_X}^\mathbf{L} L))
$$
and by the vanishing of cohomology sheaves, this is equal to
$H^0(X, \tau_{\geq 1}(K^\vee \otimes_{\mathcal{O}_X}^\mathbf{L} L))$
which is zero by Derived Categories, Lemma
\href{https://stacks.math.columbia.edu/tag/05TC}{lemma negative vanishing (Tag 05TC)}.
It follows that $\Hom_{D(\mathcal{O}_X)}(K, L)$ is zero if
$H^i(L) = 0$ for $i \in [-b - N, -a]$.
\end{proof}


\begin{lemma}
\label{lemma-lift-quasi-coherent}
Let $X$ be a scheme and let $j : U \to X$ be a quasi-compact
open immersion. The functors
$$
D_\QCoh(\mathcal{O}_X) \to D_\QCoh(\mathcal{O}_U)
\quad\text{and}\quad
D^+_\QCoh(\mathcal{O}_X) \to D^+_\QCoh(\mathcal{O}_U)
$$
are essentially surjective. If $X$ is quasi-compact, then the functors
$$
D^-_\QCoh(\mathcal{O}_X) \to D^-_\QCoh(\mathcal{O}_U)
\quad\text{and}\quad
D^b_\QCoh(\mathcal{O}_X) \to D^b_\QCoh(\mathcal{O}_U)
$$
are essentially surjective.
\end{lemma}

\begin{proof}
The argument preceding the lemma applies for the first case because $Rj_*$
maps $D_\QCoh(\mathcal{O}_U)$ into $D_\QCoh(\mathcal{O}_X)$
by Lemma \ref{lemma-quasi-coherence-direct-image}.
It is clear that $Rj_*$ maps
$D^+_\QCoh(\mathcal{O}_U)$ into
$D^+_\QCoh(\mathcal{O}_X)$
which implies the statement on bounded below complexes.
Finally, Lemma \ref{lemma-quasi-coherence-direct-image}
guarantees that $Rj_*$ maps
$D^-_\QCoh(\mathcal{O}_U)$ into
$D^-_\QCoh(\mathcal{O}_X)$
if $X$ is quasi-compact. Combining these two we obtain the last statement.
\end{proof}


\begin{lemma}
\label{lemma-quasi-coherence-direct-image}
Let $f : X \to S$ be a morphism of schemes.
Assume that $f$ is quasi-separated and quasi-compact.
\begin{enumerate}
\item The functor $Rf_*$ sends $D_\QCoh(\mathcal{O}_X)$
into $D_\QCoh(\mathcal{O}_S)$.
\item If $S$ is quasi-compact, there exists an integer $N = N(X, S, f)$
such that for an object $E$ of $D_\QCoh(\mathcal{O}_X)$
with $H^m(E) = 0$ for $m > 0$ we have
$H^m(Rf_*E) = 0$ for $m \geq N$.
\item In fact, if $S$ is quasi-compact we can find $N = N(X, S, f)$
such that for every morphism of schemes $S' \to S$
the same conclusion holds for the functor $R(f')_*$
where $f' : X' \to S'$ is the base change of $f$.
\end{enumerate}
\end{lemma}

\begin{proof}
Let $E$ be an object of $D_\QCoh(\mathcal{O}_X)$. To prove (1) we have to
show that $Rf_*E$ has quasi-coherent cohomology sheaves. The question is local
on $S$, hence we may assume $S$ is quasi-compact. Pick $N = N(X, S, f)$ as in
Cohomology of Schemes, Lemma
\href{https://stacks.math.columbia.edu/tag/01XJ}{lemma quasi coherence higher direct images (Tag 01XJ)}.
Thus $R^pf_*\mathcal{F} = 0$ for all quasi-coherent $\mathcal{O}_X$-modules
$\mathcal{F}$ and all $p \geq N$ and the same remains true after base change.

\medskip\noindent
First, assume $E$ is bounded below. We will show (1) and (2) and (3) hold
for such $E$ with our choice of $N$. In this case we can for example use the
spectral sequence
$$
R^pf_*H^q(E) \Rightarrow R^{p + q}f_*E
$$
(Derived Categories, Lemma \href{https://stacks.math.columbia.edu/tag/015J}{lemma two ss complex functor (Tag 015J)}),
the quasi-coherence of $R^pf_*H^q(E)$, and the vanishing of $R^pf_*H^q(E)$
for $p \geq N$ to see that (1), (2), and (3) hold in this case.

\medskip\noindent
Next we prove (2) and (3). Say $H^m(E) = 0$ for $m > 0$.
Let $U \subset S$ be affine open. By Cohomology of Schemes, Lemma
\href{https://stacks.math.columbia.edu/tag/01XK}{lemma quasi coherence higher direct images application (Tag 01XK)}
and our choice of $N$
we have $H^p(f^{-1}(U), \mathcal{F}) = 0$ for $p \geq N$
and any quasi-coherent $\mathcal{O}_X$-module $\mathcal{F}$.
Hence we may apply Lemma \ref{lemma-application-nice-K-injective}
to the functor $\Gamma(f^{-1}(U), -)$ to see that
$$
R\Gamma(U, Rf_*E) = R\Gamma(f^{-1}(U), E)
$$
has vanishing cohomology in degrees $\geq N$. Since this holds for
all $U \subset S$ affine open we conclude that $H^m(Rf_*E) = 0$
for $m \geq N$.

\medskip\noindent
Next, we prove (1) in the general case. Recall that there is a
distinguished triangle
$$
\tau_{\leq -n - 1}E \to E \to \tau_{\geq -n}E \to
(\tau_{\leq -n - 1}E)[1]
$$
in $D(\mathcal{O}_X)$, see Derived Categories, Remark
\href{https://stacks.math.columbia.edu/tag/08J5}{remark truncation distinguished triangle (Tag 08J5)}.
By (2) we see that $Rf_*\tau_{\leq -n - 1}E$
has vanishing cohomology sheaves in degrees $\geq -n + N$.
Thus, given an integer $q$ we see that $R^qf_*E$ is equal
to $R^qf_*\tau_{\geq -n}E$ for some $n$ and the result
above applies.
\end{proof}


\begin{lemma}
\label{lemma-affine-compare-bounded}
Let $X = \Spec(A)$ be an affine scheme. All the functors in the diagram
$$
\xymatrix{
D(\QCoh(\mathcal{O}_X)) \ar[rr]_{(\href{https://stacks.math.columbia.edu/tag/06VT}{equation compare (Tag 06VT)})}
& &
D_\QCoh(\mathcal{O}_X) \ar[ld]^{R\Gamma(X, -)} \\
& D(A) \ar[lu]^{\widetilde{\ \ }}
}
$$
are equivalences of triangulated categories. Moreover, for $E$ in
$D_\QCoh(\mathcal{O}_X)$ we have $H^0(X, E) = H^0(X, H^0(E))$.
\end{lemma}

\begin{proof}
The functor $R\Gamma(X, -)$ gives a functor
$D(\mathcal{O}_X) \to D(A)$ and hence by restriction a functor
\begin{equation}
\label{equation-back}
R\Gamma(X, -) : D_\QCoh(\mathcal{O}_X) \longrightarrow D(A).
\end{equation}
We will show this functor is quasi-inverse to (\href{https://stacks.math.columbia.edu/tag/06VT}{equation compare (Tag 06VT)})
via the equivalence between quasi-coherent modules on $X$ and
the category of $A$-modules.

\medskip\noindent
Elucidation. Denote $(Y, \mathcal{O}_Y)$ the one point space with sheaf
of rings given by $A$. Denote
$\pi : (X, \mathcal{O}_X) \to (Y, \mathcal{O}_Y)$
the obvious morphism of ringed spaces.
Then $R\Gamma(X, -)$ can be identified with $R\pi_*$ and the functor
(\href{https://stacks.math.columbia.edu/tag/06VT}{equation compare (Tag 06VT)}) via the equivalence
$\textit{Mod}(\mathcal{O}_Y) = \text{Mod}_A = \QCoh(\mathcal{O}_X)$
can be identified with $L\pi^* = \pi^* = \widetilde{\ }$ (see
Modules, Lemma \href{https://stacks.math.columbia.edu/tag/01BH}{lemma construct quasi coherent sheaves (Tag 01BH)} and
Schemes, Lemmas \href{https://stacks.math.columbia.edu/tag/01I7}{lemma compare constructions (Tag 01I7)} and
\href{https://stacks.math.columbia.edu/tag/01IB}{lemma equivalence quasi coherent (Tag 01IB)}). Thus the functors
$$
\xymatrix{
D(A) \ar@<1ex>[r] & D(\mathcal{O}_X) \ar@<1ex>[l]
}
$$
are adjoint (by Cohomology, Lemma \href{https://stacks.math.columbia.edu/tag/079W}{lemma adjoint (Tag 079W)}). In
particular we obtain canonical adjunction mappings
$$
a : \widetilde{R\Gamma(X, E)} \longrightarrow E
$$
for $E$ in $D(\mathcal{O}_X)$ and
$$
b : M^\bullet \longrightarrow R\Gamma(X, \widetilde{M^\bullet})
$$
for $M^\bullet$ a complex of $A$-modules.

\medskip\noindent
Let $E$ be an object of $D_\QCoh(\mathcal{O}_X)$. We may apply
Lemma \ref{lemma-application-nice-K-injective}
to the functor $F(-) = \Gamma(X, -)$
with $N = 1$ by Cohomology of Schemes, Lemma
\href{https://stacks.math.columbia.edu/tag/01XB}{lemma quasi coherent affine cohomology zero (Tag 01XB)}.
Hence
$$
H^0(R\Gamma(X, E)) = H^0(R\Gamma(X, \tau_{\geq 0}E)) = \Gamma(X, H^0(E))
$$
(the last equality by definition of the canonical truncation).
Using this we will show that the adjunction mappings $a$ and $b$
induce isomorphisms $H^0(a)$ and $H^0(b)$. Thus $a$ and $b$
are quasi-isomorphisms (as the statement is invariant under shifts)
and the lemma is proved.

\medskip\noindent
In both cases we use that $\widetilde{\ }$ is an exact functor
(Schemes, Lemma \href{https://stacks.math.columbia.edu/tag/01HV}{lemma spec sheaves (Tag 01HV)}). Namely, this
implies that
$$
H^0\left(\widetilde{R\Gamma(X, E)}\right) =
\widetilde{H^0(R\Gamma(X, E))} =
\widetilde{\Gamma(X, H^0(E))}
$$
which is equal to $H^0(E)$ because $H^0(E)$ is quasi-coherent. Thus
$H^0(a)$ is an isomorphism. For the other direction we have
$$
H^0(R\Gamma(X, \widetilde{M^\bullet})) =
\Gamma(X, H^0(\widetilde{M^\bullet})) =
\Gamma(X, \widetilde{H^0(M^\bullet)}) =
H^0(M^\bullet)
$$
which proves that $H^0(b)$ is an isomorphism.
\end{proof}


\begin{lemma}
\label{lemma-pseudo-coherent}
Let $X$ be a scheme. If $E$ is an $m$-pseudo-coherent
object of $D(\mathcal{O}_X)$, then $H^i(E)$ is a quasi-coherent
$\mathcal{O}_X$-module for $i > m$ and $H^m(E)$ is a quotient
of a quasi-coherent $\mathcal{O}_X$-module.
If $E$ is pseudo-coherent, then $E$ is an object of
$D_\QCoh(\mathcal{O}_X)$.
\end{lemma}

\begin{proof}
Locally on $X$ there exists a strictly perfect complex $\mathcal{E}^\bullet$
such that $H^i(E)$ is isomorphic to $H^i(\mathcal{E}^\bullet)$ for $i > m$
and $H^m(E)$ is a quotient of $H^m(\mathcal{E}^\bullet)$. The sheaves
$\mathcal{E}^i$ are direct summands of finite free modules,
hence quasi-coherent. The lemma follows.
\end{proof}


\begin{lemma}
\label{lemma-pseudo-coherent-affine}
Let $X = \Spec(A)$ be an affine scheme. Let $M^\bullet$ be a
complex of $A$-modules and let $E$ be the corresponding object
of $D(\mathcal{O}_X)$. Then $E$ is an $m$-pseudo-coherent
(resp.\ pseudo-coherent) as an object of $D(\mathcal{O}_X)$
if and only if $M^\bullet$ is $m$-pseudo-coherent (resp.\ pseudo-coherent)
as a complex of $A$-modules.
\end{lemma}

\begin{proof}
It is immediate from the definitions that if $M^\bullet$ is
$m$-pseudo-coherent, so is $E$. To prove the converse, assume
$E$ is $m$-pseudo-coherent. As $X = \Spec(A)$ is quasi-compact with
a basis for the topology given by standard opens, we can find a standard
open covering $X = D(f_1) \cup \ldots \cup D(f_n)$ and strictly
perfect complexes $\mathcal{E}_i^\bullet$ on $D(f_i)$ and
maps $\alpha_i : \mathcal{E}_i^\bullet \to E|_{U_i}$ inducing
isomorphisms on $H^j$ for $j > m$ and surjections on $H^m$.
By Cohomology, Lemma \href{https://stacks.math.columbia.edu/tag/08C9}{lemma local actual (Tag 08C9)}
after refining the open covering
we may assume $\alpha_i$ is given by a map of complexes
$\mathcal{E}_i^\bullet \to \widetilde{M^\bullet}|_{U_i}$
for each $i$. By Modules, Lemma
\href{https://stacks.math.columbia.edu/tag/0BCI}{lemma direct summand of locally free is locally free (Tag 0BCI)}
the terms $\mathcal{E}_i^n$ are finite locally free modules.
Hence after refining the open covering we may assume each
$\mathcal{E}_i^n$ is a finite free $\mathcal{O}_{U_i}$-module.
From the definition it follows that $M^\bullet_{f_i}$ is
an $m$-pseudo-coherent complex of $A_{f_i}$-modules.
We conclude by applying
More on Algebra, Lemma \ref{more-algebra-lemma-glue-pseudo-coherent}.

\medskip\noindent
The case ``pseudo-coherent'' follows from the fact that $E$ is
pseudo-coherent if and only if $E$ is $m$-pseudo-coherent for
all $m$ (by definition) and the same is true for $M^\bullet$
by More on Algebra, Lemma \href{https://stacks.math.columbia.edu/tag/064U}{lemma pseudo coherent (Tag 064U)}.
\end{proof}


\begin{lemma}
\label{more-algebra-lemma-cone-pseudo-coherent}
Let $R$ be a ring and $m \in \mathbf{Z}$.
Let $(K^\bullet, L^\bullet, M^\bullet, f, g, h)$ be a distinguished
triangle in $D(R)$.
\begin{enumerate}
\item If $K^\bullet$ is $(m + 1)$-pseudo-coherent and
$L^\bullet$ is $m$-pseudo-coherent then $M^\bullet$ is
$m$-pseudo-coherent.
\item If $K^\bullet, M^\bullet$ are $m$-pseudo-coherent, then
$L^\bullet$ is $m$-pseudo-coherent.
\item If $L^\bullet$ is $(m + 1)$-pseudo-coherent and $M^\bullet$
is $m$-pseudo-coherent, then $K^\bullet$ is $(m + 1)$-pseudo-coherent.
\end{enumerate}
\end{lemma}

\begin{proof}
Proof of (1). Choose $\alpha : P^\bullet \to K^\bullet$
with $P^\bullet$ a bounded complex of finite free modules
such that $H^i(\alpha)$ is an isomorphism for $i > m + 1$ and
surjective for $i = m + 1$. We may replace $P^\bullet$ by
$\sigma_{\geq m + 1}P^\bullet$ and hence we may assume that $P^i = 0$
for $i < m + 1$. Choose $\beta : E^\bullet \to L^\bullet$ with $E^\bullet$
a bounded complex of finite free modules such that
$H^i(\beta)$ is an isomorphism for $i > m$ and
surjective for $i = m$. By
Derived Categories,
Lemma \href{https://stacks.math.columbia.edu/tag/064E}{lemma lift map (Tag 064E)}
we can find a map $\gamma : P^\bullet \to E^\bullet$ such that the diagram
$$
\xymatrix{
K^\bullet \ar[r] & L^\bullet \\
P^\bullet \ar[u] \ar[r]^\gamma & E^\bullet \ar[u]_\beta
}
$$
is commutative in $D(R)$. The cone $C(\gamma)^\bullet$ is a bounded
complex of finite free $R$-modules, and the commutativity of the
diagram implies that there exists a morphism of distinguished triangles
$$
(P^\bullet, E^\bullet, C(\gamma)^\bullet)
\longrightarrow
(K^\bullet, L^\bullet, M^\bullet).
$$
It follows from the induced map on long exact cohomology sequences and
Homology, Lemmas \href{https://stacks.math.columbia.edu/tag/05QA}{lemma four lemma (Tag 05QA)} and
\href{https://stacks.math.columbia.edu/tag/05QB}{lemma five lemma (Tag 05QB)}
that $C(\gamma)^\bullet \to M^\bullet$ induces an isomorphism
on cohomology in degrees $> m$ and a surjection in degree $m$.
Hence $M^\bullet$ is $m$-pseudo-coherent.

\medskip\noindent
Assertions (2) and (3) follow from (1) by rotating the distinguished
triangle.
\end{proof}


\begin{lemma}
\label{more-algebra-lemma-finite-cohomology}
Let $R$ be a ring. Let $K^\bullet$ be a complex of $R$-modules.
Let $m \in \mathbf{Z}$.
\begin{enumerate}
\item If $K^\bullet$ is $m$-pseudo-coherent and $H^i(K^\bullet) = 0$
for $i > m$, then $H^m(K^\bullet)$ is a finite type $R$-module.
\item If $K^\bullet$ is $m$-pseudo-coherent and $H^i(K^\bullet) = 0$
for $i > m + 1$, then $H^{m + 1}(K^\bullet)$ is a finitely presented
$R$-module.
\end{enumerate}
\end{lemma}

\begin{proof}
Proof of (1). Choose a bounded complex $E^\bullet$ of finite projective
$R$-modules and a map $\alpha : E^\bullet \to K^\bullet$ which induces
an isomorphism on cohomology in degrees $> m$ and a surjection in degree $m$.
It is clear that it suffices to prove the result for $E^\bullet$.
Let $n$ be the largest integer such that $E^n \not = 0$.
If $n = m$, then the result is clear.
If $n > m$, then $E^{n - 1} \to E^n$ is surjective as
$H^n(E^\bullet) = 0$. As $E^n$ is finite projective we see that
$E^{n - 1} = E' \oplus E^n$. Hence it suffices to prove the result
for the complex $(E')^\bullet$ which is the same as $E^\bullet$
except has $E'$ in degree $n - 1$ and $0$ in degree $n$.
We win by induction on $n$.

\medskip\noindent
Proof of (2). Choose a bounded complex $E^\bullet$ of finite projective
$R$-modules and a map $\alpha : E^\bullet \to K^\bullet$ which induces
an isomorphism on cohomology in degrees $> m$ and a surjection in degree $m$.
As in the proof of (1) we can reduce to the case that $E^i = 0$ for
$i > m + 1$. Then we see that
$H^{m + 1}(K^\bullet) \cong
H^{m + 1}(E^\bullet) = \Coker(E^m \to E^{m + 1})$
which is of finite presentation.
\end{proof}


\begin{lemma}
\label{more-algebra-lemma-glue-pseudo-coherent}
Let $R$ be a ring. Let $f_1, \ldots, f_r \in R$ be elements which
generate the unit ideal. Let $m \in \mathbf{Z}$. Let $K^\bullet$
be a complex of $R$-modules. If for each $i$ the complex
$K^\bullet \otimes_R R_{f_i}$ is $m$-pseudo-coherent
(resp.\ pseudo-coherent), then $K^\bullet$ is $m$-pseudo-coherent
(resp.\ pseudo-coherent).
\end{lemma}

\begin{proof}
We will use without further mention that $- \otimes_R R_{f_i}$ is
an exact functor and that therefore
$$
H^i(K^\bullet)_{f_i} =
H^i(K^\bullet) \otimes_R R_{f_i} = H^i(K^\bullet \otimes_R R_{f_i}).
$$
Assume $K^\bullet \otimes_R R_{f_i}$ is $m$-pseudo-coherent
for $i = 1, \ldots, r$. Let $n \in \mathbf{Z}$ be the largest
integer such that $H^n(K^\bullet \otimes_R R_{f_i})$ is nonzero
for some $i$. This implies in particular that $H^i(K^\bullet) = 0$
for $i > n$ (and that $H^n(K^\bullet) \not = 0$) see
Algebra, Lemma \ref{algebra-lemma-cover}.
We will prove the lemma by induction on $n - m$.
If $n < m$, then the lemma is true by
Lemma \href{https://stacks.math.columbia.edu/tag/064W}{lemma recognize pseudo coherent (Tag 064W)}.
If $n \geq m$, then $H^n(K^\bullet)_{f_i}$ is a finite $R_{f_i}$-module
for each $i$, see
Lemma \ref{more-algebra-lemma-finite-cohomology}.
Hence $H^n(K^\bullet)$ is a finite $R$-module, see
Algebra, Lemma \ref{algebra-lemma-cover}.
Choose a finite free $R$-module $E$ and a surjection $E \to H^n(K^\bullet)$.
As $E$ is projective we can lift this to a map of complexes
$\alpha : E[-n] \to K^\bullet$. Then the cone $C(\alpha)^\bullet$ has
vanishing cohomology in degrees $\geq n$. On the other hand, the
complexes $C(\alpha)^\bullet \otimes_R R_{f_i}$ are $m$-pseudo-coherent
for each $i$, see
Lemma \ref{more-algebra-lemma-cone-pseudo-coherent}.
Hence by induction we see that $C(\alpha)^\bullet$ is $m$-pseudo-coherent
as a complex of $R$-modules. Applying
Lemma \ref{more-algebra-lemma-cone-pseudo-coherent}
once more we conclude.
\end{proof}


\begin{proposition}
\label{algebra-proposition-localization-exact}
\begin{slogan}
Localization is exact.
\end{slogan}
Let $L\xrightarrow{u} M\xrightarrow{v} N$ be an exact sequence
of $R$-modules. Then
$S^{-1}L \to S^{-1}M \to S^{-1}N$ is also exact.
\end{proposition}

\begin{proof}
First it is clear that $S^{-1}L \to S^{-1}M \to S^{-1}N$ is a complex
since localization is a functor. Next suppose that $x/s$ maps to zero
in $S^{-1}N$ for some $x/s \in S^{-1}M$. Then by definition there is a
$t\in S$ such that $v(xt) = v(x)t = 0$ in $N$, which means
$xt \in \Ker(v)$. By the exactness of $L \to M \to N$ we have
$xt = u(y)$ for some $y$ in $L$. Then $x/s$ is the image of $y/st$.
This proves the exactness.
\end{proof}


\begin{lemma}
\label{algebra-lemma-cover}
\begin{slogan}
Zariski-local properties of modules and algebras
\end{slogan}
Let $R$ be a ring. Let $M$ be an $R$-module. Let $S$ be an $R$-algebra.
Suppose that $f_1, \ldots, f_n$ is a finite list of
elements of $R$ such that $\bigcup D(f_i) = \Spec(R)$,
in other words $(f_1, \ldots, f_n) = R$.
\begin{enumerate}
\item If each $M_{f_i} = 0$ then $M = 0$.
\item If each $M_{f_i}$ is a finite $R_{f_i}$-module,
then $M$ is a finite $R$-module.
\item If each $M_{f_i}$ is a finitely presented $R_{f_i}$-module,
then $M$ is a finitely presented $R$-module.
\item Let $M \to N$ be a map of $R$-modules. If $M_{f_i} \to N_{f_i}$
is an isomorphism for each $i$ then $M \to N$ is an isomorphism.
\item Let $0 \to M'' \to M \to M' \to 0$ be a complex of $R$-modules.
If $0 \to M''_{f_i} \to M_{f_i} \to M'_{f_i} \to 0$ is exact for each $i$,
then $0 \to M'' \to M \to M' \to 0$ is exact.
\item If each $R_{f_i}$ is Noetherian, then $R$ is Noetherian.
\item If each $S_{f_i}$ is a finite type $R_{f_i}$-algebra, then
$S$ is a finite type $R$-algebra.
\item If each $S_{f_i}$ is of finite presentation over $R_{f_i}$, then
$S$ is a finitely presented $R$-algebra.
\end{enumerate}
\end{lemma}

\begin{proof}
We prove each of the parts in turn.
\begin{enumerate}
\item By Proposition \href{https://stacks.math.columbia.edu/tag/02C6}{proposition localize twice (Tag 02C6)}
this implies $M_\mathfrak p = 0$ for all $\mathfrak p \in \Spec(R)$,
so we conclude by Lemma \href{https://stacks.math.columbia.edu/tag/00HN}{lemma characterize zero local (Tag 00HN)}.
\item For each $i$ take a finite generating set $X_i$ of $M_{f_i}$.
Without loss of generality, we may assume that the elements of $X_i$ are
in the image of the localization map $M \rightarrow M_{f_i}$, so we take
a finite set $Y_i$ of preimages of the elements of $X_i$ in $M$. Let $Y$
be the union of these sets. This is still a finite set.
Consider the obvious $R$-linear map $R^Y \rightarrow M$ sending the basis
element $e_y$ to $y$. By assumption this map is surjective after localizing
at an arbitrary prime ideal $\mathfrak p$ of $R$, so it is surjective by
Lemma \href{https://stacks.math.columbia.edu/tag/00HN}{lemma characterize zero local (Tag 00HN)}
and $M$ is finitely generated.
\item By (2) we have a short exact sequence
$$
0 \rightarrow K \rightarrow R^n \rightarrow M \rightarrow 0
$$
Since localization is an exact functor and $M_{f_i}$ is finitely
presented we see that $K_{f_i}$ is finitely generated for all
$1 \leq i \leq n$ by Lemma \href{https://stacks.math.columbia.edu/tag/0519}{lemma extension (Tag 0519)}.
By (2) this implies that $K$ is a finite $R$-module and therefore
$M$ is finitely presented.
\item By Proposition \href{https://stacks.math.columbia.edu/tag/02C6}{proposition localize twice (Tag 02C6)}
the assumption implies that the induced morphism
on localizations at all prime ideals is an isomorphism, so we conclude
by Lemma \href{https://stacks.math.columbia.edu/tag/00HN}{lemma characterize zero local (Tag 00HN)}.
\item By Proposition \href{https://stacks.math.columbia.edu/tag/02C6}{proposition localize twice (Tag 02C6)} the assumption
implies that the induced
sequence of localizations at all prime ideals is short exact, so we
conclude by Lemma \href{https://stacks.math.columbia.edu/tag/00HN}{lemma characterize zero local (Tag 00HN)}.
\item We will show that every ideal of $R$ has a finite generating set:
For this, let $I \subset R$ be an arbitrary ideal. By
Proposition \ref{algebra-proposition-localization-exact}
each $I_{f_i} \subset R_{f_i}$ is an ideal. These are all
finitely generated by assumption, so we conclude by (2).
\item For each $i$ take a finite generating set $X_i$ of $S_{f_i}$.
Without loss of generality, we may assume that the elements of $X_i$
are in the image of the localization map $S \rightarrow S_{f_i}$, so
we take a finite set $Y_i$ of preimages of the elements of $X_i$ in
$S$. Let $Y$ be the union of these sets. This is still a finite set.
Consider the algebra homomorphism $R[X_y]_{y \in Y} \rightarrow S$
induced by $Y$. Since it is an algebra homomorphism, the image $T$
is an $R$-submodule of the $R$-module $S$, so we can consider the
quotient module $S/T$. By assumption, this is zero if we localize
at the $f_i$, so it is zero by (1) and therefore $S$ is an
$R$-algebra of finite type.
\item By the previous item, there exists a surjective $R$-algebra
homomorphism $R[X_1, \ldots, X_n] \rightarrow S$. Let $K$ be the kernel
of this map. This is an ideal in $R[X_1, \ldots, X_n]$, finitely generated
in each localization at $f_i$. Since the $f_i$ generate the unit ideal
in $R$, they also generate the unit ideal in $R[X_1, \ldots, X_n]$, so an
application of (2) finishes the proof.
\end{enumerate}
\end{proof}


\begin{lemma}
\label{cohomology-lemma-last-one-flat}
Let $(X, \mathcal{O}_X)$ be a ringed space.
Let $\mathcal{E}^\bullet$ be a bounded above complex of flat
$\mathcal{O}_X$-modules with tor-amplitude in $[a, b]$.
Then $\Coker(d_{\mathcal{E}^\bullet}^{a - 1})$ is a flat
$\mathcal{O}_X$-module.
\end{lemma}

\begin{proof}
As $\mathcal{E}^\bullet$ is a bounded above complex of flat modules we see that
$\mathcal{E}^\bullet \otimes_{\mathcal{O}_X} \mathcal{F} =
\mathcal{E}^\bullet \otimes_{\mathcal{O}_X}^{\mathbf{L}} \mathcal{F}$
for any $\mathcal{O}_X$-module $\mathcal{F}$.
Hence for every $\mathcal{O}_X$-module $\mathcal{F}$ the sequence
$$
\mathcal{E}^{a - 2} \otimes_{\mathcal{O}_X} \mathcal{F} \to
\mathcal{E}^{a - 1} \otimes_{\mathcal{O}_X} \mathcal{F} \to
\mathcal{E}^a \otimes_{\mathcal{O}_X} \mathcal{F}
$$
is exact in the middle. Since
$\mathcal{E}^{a - 2} \to \mathcal{E}^{a - 1} \to \mathcal{E}^a \to
\Coker(d^{a - 1}) \to 0$
is a flat resolution this implies that
$\text{Tor}_1^{\mathcal{O}_X}(\Coker(d^{a - 1}), \mathcal{F}) = 0$
for all $\mathcal{O}_X$-modules $\mathcal{F}$. This means that
$\Coker(d^{a - 1})$ is flat, see Lemma \href{https://stacks.math.columbia.edu/tag/08BQ}{lemma flat tor zero (Tag 08BQ)}.
\end{proof}


\begin{lemma}
\label{cohomology-lemma-perfect}
Let $(X, \mathcal{O}_X)$ be a ringed space.
Let $E$ be an object of $D(\mathcal{O}_X)$.
The following are equivalent
\begin{enumerate}
\item $E$ is perfect, and
\item $E$ is pseudo-coherent and locally has finite tor dimension.
\end{enumerate}
\end{lemma}

\begin{proof}
Assume (1). By definition this means there exists an open covering
$X = \bigcup U_i$ such that $E|_{U_i}$ is represented by a
strictly perfect complex. Thus $E$ is pseudo-coherent (i.e.,
$m$-pseudo-coherent for all $m$) by
Lemma \href{https://stacks.math.columbia.edu/tag/08CC}{lemma pseudo coherent independent representative (Tag 08CC)}.
Moreover, a direct summand of a finite free module is flat, hence
$E|_{U_i}$ has finite Tor dimension by
Lemma \href{https://stacks.math.columbia.edu/tag/08CI}{lemma tor amplitude (Tag 08CI)}. Thus (2) holds.

\medskip\noindent
Assume (2). After replacing $X$ by the members of an open covering
we may assume there exist integers $a \leq b$ such that $E$
has tor amplitude in $[a, b]$. Since $E$ is $m$-pseudo-coherent
for all $m$ we conclude using Lemma \href{https://stacks.math.columbia.edu/tag/08CP}{lemma perfect precise (Tag 08CP)}.
\end{proof}


\begin{lemma}
\label{properties-lemma-finite-locally-free}
Let $X$ be a scheme.
Let $\mathcal{F}$ be a quasi-coherent $\mathcal{O}_X$-module.
The following are equivalent:
\begin{enumerate}
\item $\mathcal{F}$ is a flat $\mathcal{O}_X$-module of finite presentation,
\item $\mathcal{F}$ is $\mathcal{O}_X$-module of finite presentation and
for all $x \in X$ the stalk $\mathcal{F}_x$ is a free
$\mathcal{O}_{X, x}$-module,
\item $\mathcal{F}$ is a locally free, finite type $\mathcal{O}_X$-module,
\item $\mathcal{F}$ is a finite locally free $\mathcal{O}_X$-module, and
\item $\mathcal{F}$ is an $\mathcal{O}_X$-module of finite type,
for every $x \in X$ the stalk $\mathcal{F}_x$ is a free
$\mathcal{O}_{X, x}$-module, and the function
$$
\rho_\mathcal{F} : X \to \mathbf{Z}, \quad
x \longmapsto
\dim_{\kappa(x)} \mathcal{F}_x \otimes_{\mathcal{O}_{X, x}} \kappa(x)
$$
is locally constant in the Zariski topology on $X$.
\end{enumerate}
\end{lemma}

\begin{proof}
This lemma immediately reduces to the affine case.
In this case the lemma is a reformulation of
Algebra, Lemma \href{https://stacks.math.columbia.edu/tag/00NX}{lemma finite projective (Tag 00NX)}.
The translation uses
Lemmas \href{https://stacks.math.columbia.edu/tag/01PB}{lemma finite type module (Tag 01PB)},
\href{https://stacks.math.columbia.edu/tag/01PC}{lemma finite presentation module (Tag 01PC)},
\href{https://stacks.math.columbia.edu/tag/05P0}{lemma flat module (Tag 05P0)}, and
\href{https://stacks.math.columbia.edu/tag/05JM}{lemma locally free module (Tag 05JM)}.
\end{proof}


\begin{lemma}
\label{lemma-application-nice-K-injective}
Let $X$ be a scheme. Let $F : \textit{Mod}(\mathcal{O}_X) \to \textit{Ab}$
be an additive functor and $N \geq 0$ an integer. Assume that
\begin{enumerate}
\item $F$ commutes with countable direct products,
\item $R^pF(\mathcal{F}) = 0$ for all $p \geq N$ and $\mathcal{F}$
quasi-coherent.
\end{enumerate}
Then for $E \in D_\QCoh(\mathcal{O}_X)$
\begin{enumerate}
\item $H^i(RF(\tau_{\leq a}E)) \to H^i(RF(E))$ is an isomorphism
for $i \leq a$,
\item $H^i(RF(E)) \to H^i(RF(\tau_{\geq b - N + 1}E))$ is an isomorphism
for $i \geq b$,
\item if $H^i(E) = 0$ for $i \not \in [a, b]$ for some
$-\infty \leq a \leq b \leq \infty$, then $H^i(RF(E)) = 0$
for $i \not \in [a, b + N - 1]$.
\end{enumerate}
\end{lemma}

\begin{proof}
Statement (1) is
Derived Categories, Lemma \href{https://stacks.math.columbia.edu/tag/05TC}{lemma negative vanishing (Tag 05TC)}.

\medskip\noindent
Proof of statement (2). Write $E_n = \tau_{\geq -n}E$. We have
$E = R\lim E_n$, see Lemma \href{https://stacks.math.columbia.edu/tag/08D3}{lemma nice K injective (Tag 08D3)}. Thus
$RF(E) = R\lim RF(E_n)$ in $D(\textit{Ab})$ by Injectives, Lemma
\href{https://stacks.math.columbia.edu/tag/08U1}{lemma RF commutes with Rlim (Tag 08U1)}. Thus for every $i \in \mathbf{Z}$
we have a short exact sequence
$$
0 \to R^1\lim H^{i - 1}(RF(E_n)) \to H^i(RF(E)) \to \lim H^i(RF(E_n)) \to 0
$$
see More on Algebra, Remark
\href{https://stacks.math.columbia.edu/tag/08U5}{remark compare derived limit (Tag 08U5)}.
To prove (2) we will show that the term on the left is zero
and that the term on the right equals $H^i(RF(E_{-b + N - 1}))$
for any $b$ with $i \geq b$.

\medskip\noindent
For every $n$ we have a distinguished triangle
$$
H^{-n}(E)[n] \to E_n \to E_{n - 1} \to H^{-n}(E)[n + 1]
$$
(Derived Categories, Remark
\href{https://stacks.math.columbia.edu/tag/08J5}{remark truncation distinguished triangle (Tag 08J5)})
in $D(\mathcal{O}_X)$. Since $H^{-n}(E)$ is quasi-coherent we have
$$
H^i(RF(H^{-n}(E)[n])) = R^{i + n}F(H^{-n}(E)) = 0
$$
for $i + n \geq N$ and
$$
H^i(RF(H^{-n}(E)[n + 1])) = R^{i + n + 1}F(H^{-n}(E)) = 0
$$
for $i + n + 1 \geq N$. We conclude that
$$
H^i(RF(E_n)) \to H^i(RF(E_{n - 1}))
$$
is an isomorphism for $n \geq N - i$. Thus the systems $H^i(RF(E_n))$ all
satisfy the ML condition and the $R^1\lim$ term in our short exact
sequence is zero (see discussion in
More on Algebra, Section \href{https://stacks.math.columbia.edu/tag/07KV}{section Rlim (Tag 07KV)}).
Moreover, the system $H^i(RF(E_n))$ is constant starting
with $n = N - i - 1$ as desired.

\medskip\noindent
Proof of (3). Under the assumption on $E$ we have
$\tau_{\leq a - 1}E = 0$ and we get the vanishing
of $H^i(RF(E))$ for $i \leq a - 1$ from (1).
Similarly, we have $\tau_{\geq b + 1}E = 0$ and hence
we get the vanishing of $H^i(RF(E))$ for $i \geq b + N$ from
part (2).
\end{proof}


\noindent
Let $U \subset X$ be an open subspace of a ringed space
and denote $j : U \to X$ the inclusion morphism. The functor
$D(\mathcal{O}_X) \to D(\mathcal{O}_U)$ is essentially surjective as
$Rj_*$ is a right inverse to restriction.
In this section we extend this to complexes with quasi-coherent cohomology
sheaves, etc.


\noindent
The notion of localization of a ring can be generalized to the
localization of a module. Let $A$ be a ring, $S$ a multiplicative
subset of $A$ and $M$ an $A$-module. We define a relation on
$M \times S$ as follows
$$
(m, s) \sim (n, t)
\Leftrightarrow
\exists u\in S \text{ such that } (mt-ns)u = 0
$$
This is clearly an equivalence relation. Denote by $m/s$ (or
$\frac{m}{s}$) be the equivalence class of $(m, s)$ and $S^{-1}M$ be
the set of all equivalence classes. Define the addition and scalar
multiplication as follows
$$
m/s + n/t = (mt + ns)/st,\quad
m/s\cdot n/t = mn/st
$$
It is clear that this makes $S^{-1}M$ an $S^{-1}A$-module.


\begin{lemma}
\label{lemma-pseudo-coherent-hocolim}
Let $X$ be a quasi-compact and quasi-separated scheme.
Let $K \in D(\mathcal{O}_X)$. The following are equivalent
\begin{enumerate}
\item $K$ is pseudo-coherent, and
\item $K = \text{hocolim} K_n$ where
$K_n$ is perfect and $\tau_{\geq -n}K_n \to \tau_{\geq -n}K$
is an isomorphism for all $n$.
\end{enumerate}
\end{lemma}

\begin{proof}
The implication (2) $\Rightarrow$ (1) is true on any ringed space.
Namely, assume (2) holds. Recall that a perfect object of the derived
category is pseudo-coherent, see
Cohomology, Lemma \ref{cohomology-lemma-perfect}.
Then it follows from the definitions that
$\tau_{\geq -n}K_n$ is $(-n + 1)$-pseudo-coherent
and hence $\tau_{\geq -n}K$ is $(-n + 1)$-pseudo-coherent,
hence $K$ is $(-n + 1)$-pseudo-coherent. This is true for
all $n$, hence $K$ is pseudo-coherent, see
Cohomology, Definition \ref{cohomology-definition-pseudo-coherent}.

\medskip\noindent
Assume (1). We start by choosing an approximation
$K_1 \to K$ of $(X, K, -2)$ by a perfect complex $K_1$, see
Definitions \ref{definition-approximation-holds} and
\ref{definition-approximation} and
Theorem \ref{theorem-approximation}.
Suppose by induction we have
$$
K_1 \to K_2 \to \ldots \to K_n \to K
$$
with $K_i$ perfect such that
such that $\tau_{\geq -i}K_i \to \tau_{\geq -i}K$ is an isomorphism
for all $1 \leq i \leq n$. Then we pick $a \leq b$ as in
Lemma \ref{lemma-ext-from-perfect-into-bounded-QCoh}
for the perfect object $K_n$. Choose an approximation
$K_{n + 1} \to K$ of $(X, K, \min(a - 1, -n - 1))$.
Choose a distinguished triangle
$$
K_{n + 1} \to K \to C \to K_{n + 1}[1]
$$
Then we see that $C \in D_\QCoh(\mathcal{O}_X)$ has
$H^i(C) = 0$ for $i \geq a$. Thus by our choice of $a, b$
we see that $\Hom_{D(\mathcal{O}_X)}(K_n, C) = 0$.
Hence the composition $K_n \to K \to C$ is zero. Hence by
Derived Categories, Lemma \href{https://stacks.math.columbia.edu/tag/0149}{lemma representable homological (Tag 0149)}
we can factor $K_n \to K$ through $K_{n + 1}$
proving the induction step.

\medskip\noindent
We still have to prove that $K = \text{hocolim} K_n$.
This follows by an application of
Derived Categories, Lemma \href{https://stacks.math.columbia.edu/tag/0CRK}{lemma cohomology of hocolim (Tag 0CRK)}
to the functors
$H^i( - ) : D(\mathcal{O}_X) \to \textit{Mod}(\mathcal{O}_X)$
and our choice of $K_n$.
\end{proof}


\begin{lemma}
\label{lemma-compute-ext-perfect}
Assumptions and notation as in Lemma \ref{lemma-ext-perfect}.
Then there are functorial isomorphisms
$$
H^i(S, K \otimes^\mathbf{L}_{\mathcal{O}_S} \mathcal{F})
\longrightarrow
\Ext^i_{\mathcal{O}_X}(E,
\mathcal{G}^\bullet \otimes_{\mathcal{O}_X} f^*\mathcal{F})
$$
for $\mathcal{F}$ quasi-coherent on $S$
compatible with boundary maps (see proof).
\end{lemma}

\begin{proof}
As in the proof of Lemma \ref{lemma-ext-perfect} let $E^\vee$ be the
dual perfect complex and recall that
$K = Rf_*(E^\vee \otimes_{\mathcal{O}_X}^\mathbf{L} \mathcal{G}^\bullet)$.
Since we also have
$$
\Ext^i_{\mathcal{O}_X}(E,
\mathcal{G}^\bullet \otimes_{\mathcal{O}_X} f^*\mathcal{F})
=
H^i(X, E^\vee \otimes^\mathbf{L}_{\mathcal{O}_X}
(\mathcal{G}^\bullet \otimes_{\mathcal{O}_X} f^*\mathcal{F}))
$$
by construction of $E^\vee$, the existence of the isomorphisms follows
from Lemma \ref{lemma-compute-tensor-perfect} applied to $E^\vee$
and $\mathcal{G}^\bullet$.
The statement on boundary maps means the following: Given a short
exact sequence $0 \to \mathcal{F}_1 \to \mathcal{F}_2 \to \mathcal{F}_3 \to 0$
then the isomorphisms fit into commutative diagrams
$$
\xymatrix{
H^i(S, K \otimes^\mathbf{L}_{\mathcal{O}_S} \mathcal{F}_3)
\ar[r] \ar[d]_\delta &
\Ext^i_{\mathcal{O}_X}(E,
\mathcal{G}^\bullet \otimes_{\mathcal{O}_X} f^*\mathcal{F}_3) \ar[d]^\delta \\
H^{i + 1}(S, K \otimes^\mathbf{L}_{\mathcal{O}_S} \mathcal{F}_1)
\ar[r] &
\Ext^{i + 1}_{\mathcal{O}_X}(E,
\mathcal{G}^\bullet \otimes_{\mathcal{O}_X} f^*\mathcal{F}_1)
}
$$
where the boundary maps come from the distinguished triangle
$$
K \otimes^\mathbf{L}_{\mathcal{O}_S} \mathcal{F}_1 \to
K \otimes^\mathbf{L}_{\mathcal{O}_S} \mathcal{F}_2 \to
K \otimes^\mathbf{L}_{\mathcal{O}_S} \mathcal{F}_3 \to
K \otimes^\mathbf{L}_{\mathcal{O}_S} \mathcal{F}_1[1]
$$
and the distinguished triangle in $D(\mathcal{O}_X)$ associated to
the short exact sequence
$$
0 \to
\mathcal{G}^\bullet \otimes_{\mathcal{O}_X} f^*\mathcal{F}_1 \to
\mathcal{G}^\bullet \otimes_{\mathcal{O}_X} f^*\mathcal{F}_2 \to
\mathcal{G}^\bullet \otimes_{\mathcal{O}_X} f^*\mathcal{F}_3 \to 0
$$
of complexes.
This sequence is exact because $\mathcal{G}$ is flat over $S$.
We omit the verification of the commutativity of the displayed diagram.
\end{proof}


\begin{lemma}
\label{lemma-ext-perfect}
Let $S$ be a Noetherian scheme. Let $f : X \to S$ be a morphism of schemes
which is locally of finite type. Let $E \in D(\mathcal{O}_X)$ be perfect.
Let $\mathcal{G}^\bullet$ be a bounded complex of coherent
$\mathcal{O}_X$-modules flat over $S$ with support proper over $S$.
Then $K = Rf_*R\SheafHom(E, \mathcal{G}^\bullet)$ is a perfect object of
$D(\mathcal{O}_S)$.
\end{lemma}

\begin{proof}
Since $E$ is a perfect complex there exists a dual perfect complex
$E^\vee$, see Cohomology, Lemma \href{https://stacks.math.columbia.edu/tag/08DQ}{lemma dual perfect complex (Tag 08DQ)}.
Observe that $R\SheafHom(E, \mathcal{G}^\bullet) =
E^\vee \otimes^\mathbf{L}_{\mathcal{O}_X} \mathcal{G}^\bullet$.
Thus the perfectness of $K$ follows from Lemma \ref{lemma-tensor-perfect}.
\end{proof}


\begin{lemma}
\label{lemma-tensor-perfect}
Let $S$ be a Noetherian scheme. Let $f : X \to S$ be a morphism of schemes
which is locally of finite type. Let $E \in D(\mathcal{O}_X)$ be perfect.
Let $\mathcal{G}^\bullet$ be a bounded complex of coherent
$\mathcal{O}_X$-modules flat over $S$ with support proper over $S$.
Then $K = Rf_*(E \otimes_{\mathcal{O}_X}^\mathbf{L} \mathcal{G}^\bullet)$
is a perfect object of $D(\mathcal{O}_S)$.
\end{lemma}

\begin{proof}
The object $K$ is perfect by Lemma \ref{lemma-perfect-direct-image}.
We check the lemma applies: Locally $E$ is isomorphic to a finite complex
of finite free $\mathcal{O}_X$-modules. Hence locally
$E \otimes^\mathbf{L}_{\mathcal{O}_X} \mathcal{G}^\bullet$ is isomorphic
to a finite complex whose terms are of the form
$$
\bigoplus\nolimits_{i = a, \ldots, b} (\mathcal{G}^i)^{\oplus r_i}
$$
for some integers $a, b, r_a, \ldots, r_b$. This immediately implies the
cohomology sheaves $H^i(E \otimes^\mathbf{L}_{\mathcal{O}_X} \mathcal{G})$
are coherent. The hypothesis on the tor dimension also follows as
$\mathcal{G}^i$ is flat over $f^{-1}\mathcal{O}_S$.
\end{proof}


\begin{lemma}
\label{lemma-compute-tensor-perfect}
Assumptions and notation as in Lemma \ref{lemma-tensor-perfect}.
Then there are functorial isomorphisms
$$
H^i(S, K \otimes^\mathbf{L}_{\mathcal{O}_S} \mathcal{F})
\longrightarrow
H^i(X, E \otimes_{\mathcal{O}_X}^\mathbf{L}
(\mathcal{G}^\bullet \otimes_{\mathcal{O}_X} f^*\mathcal{F}))
$$
for $\mathcal{F}$ quasi-coherent on $S$
compatible with boundary maps (see proof).
\end{lemma}

\begin{proof}
We have
$$
\mathcal{G}^\bullet \otimes_{\mathcal{O}_X}^\mathbf{L} Lf^*\mathcal{F} =
\mathcal{G}^\bullet \otimes_{f^{-1}\mathcal{O}_S}^\mathbf{L} f^{-1}\mathcal{F} =
\mathcal{G}^\bullet \otimes_{f^{-1}\mathcal{O}_S} f^{-1}\mathcal{F} =
\mathcal{G}^\bullet \otimes_{\mathcal{O}_X} f^*\mathcal{F}
$$
the first equality by
Cohomology, Lemma \href{https://stacks.math.columbia.edu/tag/08DE}{lemma variant derived pullback (Tag 08DE)},
the second as $\mathcal{G}^n$ is a flat $f^{-1}\mathcal{O}_S$-module, and
the third by definition of pullbacks. Hence we obtain
\begin{align*}
H^i(X, E \otimes^\mathbf{L}_{\mathcal{O}_X}
(\mathcal{G}^\bullet \otimes_{\mathcal{O}_X} f^*\mathcal{F}))
& =
H^i(X, E \otimes^\mathbf{L}_{\mathcal{O}_X} \mathcal{G}^\bullet
\otimes_{\mathcal{O}_X}^\mathbf{L} Lf^*\mathcal{F}) \\
& =
H^i(S,
Rf_*(E \otimes^\mathbf{L}_{\mathcal{O}_X} \mathcal{G}^\bullet
\otimes^\mathbf{L}_{\mathcal{O}_X} Lf^*\mathcal{F})) \\
& =
H^i(S, Rf_*(E \otimes^\mathbf{L}_{\mathcal{O}_X} \mathcal{G}^\bullet)
\otimes^\mathbf{L}_{\mathcal{O}_S} \mathcal{F}) \\
& =
H^i(S, K \otimes^\mathbf{L}_{\mathcal{O}_S} \mathcal{F}) 
\end{align*}
The first equality by the above, the second by Leray
(Cohomology, Lemma \href{https://stacks.math.columbia.edu/tag/01EZ}{lemma before Leray (Tag 01EZ)}), and
the third equality by Lemma \ref{lemma-cohomology-base-change}.
The statement on boundary maps means the following: Given a short
exact sequence $0 \to \mathcal{F}_1 \to \mathcal{F}_2 \to \mathcal{F}_3 \to 0$
of quasi-coherent $\mathcal{O}_S$-modules, the isomorphisms fit into
commutative diagrams
$$
\xymatrix{
H^i(S, K \otimes^\mathbf{L}_{\mathcal{O}_S} \mathcal{F}_3)
\ar[r] \ar[d]_\delta &
H^i(X, E \otimes^\mathbf{L}_{\mathcal{O}_X}
(\mathcal{G}^\bullet \otimes_{\mathcal{O}_X} f^*\mathcal{F}_3))
\ar[d]^\delta \\
H^{i + 1}(S, K \otimes^\mathbf{L}_{\mathcal{O}_S} \mathcal{F}_1)
\ar[r] &
H^{i + 1}(X, E \otimes^\mathbf{L}_{\mathcal{O}_X}
(\mathcal{G}^\bullet \otimes_{\mathcal{O}_X} f^*\mathcal{F}_1))
}
$$
where the boundary maps come from the distinguished triangle
$$
K \otimes^\mathbf{L}_{\mathcal{O}_S} \mathcal{F}_1 \to
K \otimes^\mathbf{L}_{\mathcal{O}_S} \mathcal{F}_2 \to
K \otimes^\mathbf{L}_{\mathcal{O}_S} \mathcal{F}_3 \to
K \otimes^\mathbf{L}_{\mathcal{O}_S} \mathcal{F}_1[1]
$$
and the distinguished triangle in $D(\mathcal{O}_X)$ associated to
the short exact sequence
$$
0 \to
\mathcal{G}^\bullet \otimes_{\mathcal{O}_X} f^*\mathcal{F}_1 \to
\mathcal{G}^\bullet \otimes_{\mathcal{O}_X} f^*\mathcal{F}_2 \to
\mathcal{G}^\bullet \otimes_{\mathcal{O}_X} f^*\mathcal{F}_3 \to 0
$$
of complexes of $\mathcal{O}_X$-modules.
This sequence is exact because $\mathcal{G}^n$ is flat over $S$.
We omit the verification of the commutativity of the displayed diagram.
\end{proof}


\begin{lemma}
\label{lemma-perfect-direct-image}
Let $S$ be a Noetherian scheme. Let $f : X \to S$ be a morphism of schemes
which is locally of finite type. Let $E \in D(\mathcal{O}_X)$ such that
\begin{enumerate}
\item $E \in D^b_{\textit{Coh}}(\mathcal{O}_X)$,
\item the support of $H^i(E)$ is proper over $S$ for all $i$, and
\item $E$ has finite tor dimension as an object of $D(f^{-1}\mathcal{O}_S)$.
\end{enumerate}
Then $Rf_*E$ is a perfect object of $D(\mathcal{O}_S)$.
\end{lemma}

\begin{proof}
By Lemma \href{https://stacks.math.columbia.edu/tag/08E2}{lemma direct image coherent (Tag 08E2)} we see that $Rf_*E$ is an object of
$D^b_{\textit{Coh}}(\mathcal{O}_S)$. Hence $Rf_*E$ is pseudo-coherent
(Lemma \href{https://stacks.math.columbia.edu/tag/08E8}{lemma identify pseudo coherent noetherian (Tag 08E8)}).
Hence it suffices to show that $Rf_*E$ has finite tor dimension, see
Cohomology, Lemma \ref{cohomology-lemma-perfect}.
By Lemma \href{https://stacks.math.columbia.edu/tag/08EA}{lemma tor qc qs (Tag 08EA)} it suffices to check that
$Rf_*(E) \otimes_{\mathcal{O}_S}^\mathbf{L} \mathcal{F}$
has universally bounded cohomology for all quasi-coherent
sheaves $\mathcal{F}$ on $S$. Bounded from above is clear as $Rf_*(E)$
is bounded from above. Let $T \subset X$ be the union of the supports
of $H^i(E)$ for all $i$. Then $T$ is proper over $S$ by assumptions
(1) and (2), see Cohomology of Schemes, Lemma
\href{https://stacks.math.columbia.edu/tag/0CYR}{lemma union closed proper over base (Tag 0CYR)}.
In particular there exists a quasi-compact open
$X' \subset X$ containing $T$. Setting $f' = f|_{X'}$ we have
$Rf_*(E) = Rf'_*(E|_{X'})$ because $E$ restricts to zero on $X \setminus T$.
Thus we may replace $X$ by $X'$ and assume $f$ is quasi-compact.
Moreover, $f$ is quasi-separated by Morphisms, Lemma
\href{https://stacks.math.columbia.edu/tag/01T7}{lemma finite type Noetherian quasi separated (Tag 01T7)}. Now
$$
Rf_*(E) \otimes_{\mathcal{O}_S}^\mathbf{L} \mathcal{F} =
Rf_*\left(E \otimes_{\mathcal{O}_X}^\mathbf{L} Lf^*\mathcal{F}\right) =
Rf_*\left(E \otimes_{f^{-1}\mathcal{O}_S}^\mathbf{L} f^{-1}\mathcal{F}\right)
$$
by
Lemma \ref{lemma-cohomology-base-change}
and
Cohomology, Lemma \href{https://stacks.math.columbia.edu/tag/08DE}{lemma variant derived pullback (Tag 08DE)}.
By assumption (3) the complex
$E \otimes_{f^{-1}\mathcal{O}_S}^\mathbf{L} f^{-1}\mathcal{F}$
has cohomology sheaves in a
given finite range, say $[a, b]$. Then $Rf_*$ of it
has cohomology in the range $[a, \infty)$ and we win.
\end{proof}


\begin{lemma}
\label{lemma-cohomology-base-change}
Let $f : X \to Y$ be a quasi-compact and quasi-separated morphism
of schemes. For $E$ in $D_\QCoh(\mathcal{O}_X)$ and
$K$ in $D_\QCoh(\mathcal{O}_Y)$ the map
$$
Rf_*(E) \otimes_{\mathcal{O}_Y}^\mathbf{L} K
\longrightarrow
Rf_*(E \otimes_{\mathcal{O}_X}^\mathbf{L} Lf^*K)
$$
defined in
Cohomology, Equation (\href{https://stacks.math.columbia.edu/tag/0B53}{equation projection formula map (Tag 0B53)})
is an isomorphism.
\end{lemma}

\begin{proof}
To check the map is an isomorphism we may work locally on $Y$.
Hence we reduce to the case that $Y$ is affine.

\medskip\noindent
Suppose that $K = \bigoplus K_i$ is a direct
sum of some complexes $K_i \in D_\QCoh(\mathcal{O}_Y)$.
If the statement holds for each $K_i$, then it holds for $K$.
Namely, the functors $Lf^*$ and $\otimes^\mathbf{L}$ preserve
direct sums by construction and $Rf_*$ commutes with direct sums
(for complexes with quasi-coherent cohomology sheaves) by
Lemma \href{https://stacks.math.columbia.edu/tag/08DZ}{lemma quasi coherence pushforward direct sums (Tag 08DZ)}.
Moreover, suppose that $K \to L \to M \to K[1]$ is a distinguished
triangle in $D_\QCoh(Y)$. Then if the statement of the
lemma holds for two of $K, L, M$, then it holds for the third
(as the functors involved are exact functors of triangulated categories).

\medskip\noindent
Assume $Y$ affine, say $Y = \Spec(A)$. The functor
$\widetilde{\ } : D(A) \to D_\QCoh(\mathcal{O}_Y)$ is an equivalence
(Lemma \ref{lemma-affine-compare-bounded}).
Let $T$ be the property for $K \in D(A)$ that
the statement of the lemma holds for $\widetilde{K}$.
The discussion above and
More on Algebra, Remark \href{https://stacks.math.columbia.edu/tag/09PB}{remark P resolution (Tag 09PB)}
shows that it suffices to prove $T$ holds for $A[k]$.
This finishes the proof, as the statement of the lemma
is clear for shifts of the structure sheaf.
\end{proof}


\begin{lemma}
\label{lemma-induction-step}
Let $X$ be a scheme. Let $X = U \cup V$ be an open covering
with $U$ quasi-compact, $V$ affine, and $U \cap V$ quasi-compact.
If approximation by perfect complexes holds on $U$,
then approximation holds on $X$.
\end{lemma}

\begin{proof}
Let $T \subset X$ be a closed subset with $X \setminus T$ retro-compact
in $X$. Let $r_U$ be the integer of Definition \ref{definition-approximation}
adapted to the pair $(U, T \cap U)$.
Set $T' = T \setminus U$. Note that
$T' \subset V$ and that $V \setminus T' = (X \setminus T) \cap U \cap V$
is quasi-compact by our assumption on $T$.
Let $r'$ be the number of affines needed to cover $V \setminus T'$.
We claim that $r = \max(r_U, r')$ works for the pair $(X, T)$.

\medskip\noindent
To see this choose a triple $(T, E, m)$ such that $E$ is
$(m - r)$-pseudo-coherent and $H^i(E)$ is supported on $T$ for
$i \geq m - r$. Let $t$ be the largest integer such that
$H^t(E)|_U$ is nonzero. (Such an integer exists as $U$ is quasi-compact
and $E|_U$ is $(m - r)$-pseudo-coherent.)
We will prove that $E$ can be approximated by induction on $t$.

\medskip\noindent
Base case: $t \leq m - r'$. This means that $H^i(E)$ is supported
on $T'$ for $i \geq m - r'$. Hence
Lemma \ref{lemma-approximation-affine}
guarantees the existence of an approximation
$P \to E|_V$ of $(T', E|_V, m)$ on $V$.
Applying Lemma \ref{lemma-open} we see that
$(T', E, m)$ can be approximated. Such an approximation
is also an approximation of $(T, E, m)$.

\medskip\noindent
Induction step. Choose an approximation $P \to E|_U$
of $(T \cap U, E|_U, m)$. This in particular gives a surjection
$H^t(P) \to H^t(E|_U)$. By
Lemma \href{https://stacks.math.columbia.edu/tag/08EK}{lemma lift perfect complex plus shift support (Tag 08EK)}
we can choose a perfect object $Q$ in $D(\mathcal{O}_V)$
supported on $T \cap V$ and an isomorphism
$Q|_{U \cap V} \to (P \oplus P[1])|_{U \cap V}$.
By Lemma \ref{lemma-lift-map} we can replace $Q$ by
$Q \otimes^\mathbf{L} I$
and assume that the map
$$
Q|_{U \cap V} \to (P \oplus P[1])|_{U \cap V}
\longrightarrow P|_{U \cap V}
\longrightarrow
E|_{U \cap V}
$$
lifts to $Q \to E|_V$. By
Cohomology, Lemma \ref{cohomology-lemma-glue}
we find an morphism $a : R \to E$ of $D(\mathcal{O}_X)$
such that $a|_U$ is isomorphic to $P \oplus P[1] \to E|_U$
and $a|_V$ isomorphic to $Q \to E|_V$.
Thus $R$ is perfect and supported on $T$
and the map $H^t(R) \to H^t(E)$ is surjective on restriction to $U$.
Choose a distinguished triangle
$$
R \to E \to E' \to R[1]
$$
Then $E'$ is $(m - r)$-pseudo-coherent
(Cohomology, Lemma \href{https://stacks.math.columbia.edu/tag/08CD}{lemma cone pseudo coherent (Tag 08CD)}),
$H^i(E')|_U = 0$ for $i \geq t$, and
$H^i(E')$ is supported on $T$ for $i \geq m - r$.
By induction we find an approximation $R' \to E'$
of $(T, E', m)$. Fit the composition $R' \to E' \to R[1]$
into a distinguished triangle $R \to R'' \to R' \to R[1]$
and extend the morphisms $R' \to E'$ and $R[1] \to R[1]$ into
a morphism of distinguished triangles
$$
\xymatrix{
R \ar[r] \ar[d] & R'' \ar[d] \ar[r] & R' \ar[d] \ar[r] & R[1] \ar[d] \\
R \ar[r] &  E \ar[r] & E' \ar[r] & R[1]
}
$$
using TR3. Then $R''$ is a perfect complex
(Cohomology, Lemma \href{https://stacks.math.columbia.edu/tag/08CR}{lemma two out of three perfect (Tag 08CR)})
supported on $T$.
An easy diagram chase shows that $R'' \to E$ is the desired
approximation.
\end{proof}


\begin{lemma}
\label{lemma-approximation-affine}
Let $X$ be an affine scheme. Then approximation holds for every
triple $(T, E, m)$ as in Definition \ref{definition-approximation-holds}
such that there exists an integer $r \geq 0$ with
\begin{enumerate}
\item $E$ is $m$-pseudo-coherent,
\item $H^i(E)$ is supported on $T$ for $i \geq m - r + 1$,
\item $X \setminus T$ is the union of $r$ affine opens.
\end{enumerate}
In particular, approximation by perfect complexes holds for affine schemes.
\end{lemma}

\begin{proof}
Say $X = \Spec(A)$. Write $T = V(f_1, \ldots, f_r)$.
(The case $r = 0$, i.e., $T = X$ follows immediately from
Lemma \ref{lemma-pseudo-coherent-affine} and the definitions.)
Let $(T, E, m)$ be a triple as in the lemma.
Let $t$ be the largest integer such that $H^t(E)$ is nonzero.
We will proceed by induction on $t$. The base case is $t < m$; in
this case the result is trivial. Now suppose that $t \geq m$. By
Cohomology, Lemma \href{https://stacks.math.columbia.edu/tag/08DN}{lemma finite cohomology (Tag 08DN)}
the sheaf $H^t(E)$ is of finite type. Since it is quasi-coherent
it is generated by finitely many sections
(Properties, Lemma \href{https://stacks.math.columbia.edu/tag/01PB}{lemma finite type module (Tag 01PB)}).
For every $s \in \Gamma(X, H^t(E)) = H^t(X, E)$
(see proof of Lemma \ref{lemma-affine-compare-bounded})
we can find an $e > 0$ and a morphism $K_e[-t] \to E$
such that $s$ is in the image of
$H^0(K_e) = H^t(K_e[-t]) \to H^t(E)$, see
Lemma \ref{lemma-represent-cohomology-class-on-closed}.
Taking a finite direct sum of these maps we obtain a map
$P \to E$ where $P$ is a perfect complex supported on $T$,
where $H^i(P) = 0$ for $i > t$, and where $H^t(P) \to E$ is
surjective. Choose a distinguished triangle
$$
P \to E \to E' \to P[1]
$$
Then $E'$ is $m$-pseudo-coherent
(Cohomology, Lemma \href{https://stacks.math.columbia.edu/tag/08CD}{lemma cone pseudo coherent (Tag 08CD)}),
$H^i(E') = 0$ for $i \geq t$, and
$H^i(E')$ is supported on $T$ for $i \geq m - r + 1$.
By induction we find an approximation $P' \to E'$
of $(T, E', m)$. Fit the composition $P' \to E' \to P[1]$
into a distinguished triangle $P \to P'' \to P' \to P[1]$
and extend the morphisms $P' \to E'$ and $P[1] \to P[1]$ into
a morphism of distinguished triangles
$$
\xymatrix{
P \ar[r] \ar[d] & P'' \ar[d] \ar[r] & P' \ar[d] \ar[r] & P[1] \ar[d] \\
P \ar[r] &  E \ar[r] & E' \ar[r] & P[1]
}
$$
using TR3. Then $P''$ is a perfect complex
(Cohomology, Lemma \href{https://stacks.math.columbia.edu/tag/08CR}{lemma two out of three perfect (Tag 08CR)})
supported on $T$.
An easy diagram chase shows that $P'' \to E$ is the desired
approximation.
\end{proof}


\begin{lemma}
\label{lemma-open}
Let $X$ be a scheme. Let $U \subset X$ be an open subscheme.
Let $(T, E, m)$ be a triple as in
Definition \ref{definition-approximation-holds}.
If
\begin{enumerate}
\item $T \subset U$,
\item approximation holds for $(T, E|_U, m)$, and
\item the sheaves $H^i(E)$ for $i \geq m$ are supported on $T$,
\end{enumerate}
then approximation holds for $(T, E, m)$.
\end{lemma}

\begin{proof}
Let $j : U \to X$ be the inclusion morphism.
If $P \to E|_U$ is an approximation of the triple $(T, E|_U, m)$
over $U$, then $j_!P = Rj_*P \to j_!(E|_U) \to E$ is an approximation
of $(T, E, m)$ over $X$.
See Cohomology, Lemmas \ref{cohomology-lemma-pushforward-restriction} and
\ref{cohomology-lemma-pushforward-perfect}.
\end{proof}


\begin{lemma}[Induction Principle]
\label{coherent-lemma-induction-principle}
\begin{reference}
\href{https://stacks.math.columbia.edu/bibliography}{Bibliography: BvdB (Proposition 3.3.1)}
\end{reference}
Let $X$ be a quasi-compact and quasi-separated scheme. Let $P$ be a property
of the quasi-compact opens of $X$. Assume that
\begin{enumerate}
\item $P$ holds for every affine open of $X$,
\item if $U$ is quasi-compact open, $V$ affine open,
$P$ holds for $U$, $V$, and $U \cap V$, then
$P$ holds for $U \cup V$.
\end{enumerate}
Then $P$ holds for every quasi-compact open of $X$
and in particular for $X$.
\end{lemma}

\begin{proof}
First we argue by induction that $P$ holds for {\it separated} quasi-compact
opens $W \subset X$. Namely, such an open can be written as
$W = U_1 \cup \ldots \cup U_n$ and we can do induction on $n$ using
property (2) with $U = U_1 \cup \ldots \cup U_{n - 1}$ and $V = U_n$.
This is allowed because
$U \cap V = (U_1 \cap U_n) \cup \ldots \cup (U_{n - 1} \cap U_n)$
is also a union of $n - 1$ affine open subschemes by
Schemes, Lemma \href{https://stacks.math.columbia.edu/tag/01KP}{lemma characterize separated (Tag 01KP)}
applied to the affine opens $U_i$ and $U_n$ of $W$.
Having said this, for any quasi-compact open $W \subset X$ we can
do induction on the number of affine opens needed to cover $W$
using the same trick as before and using that the quasi-compact open
$U_i \cap U_n$ is separated as an open subscheme of the affine scheme $U_n$.
\end{proof}


\begin{lemma}
\label{lemma-lift-map}
Let $X$ be an affine scheme. Let $U \subset X$ be a quasi-compact open.
Let $E, E'$ be objects of $D_\QCoh(\mathcal{O}_X)$ with $E$ perfect.
For every map $\alpha : E|_U \to E'|_U$ there exist maps
$$
E \xleftarrow{\beta} E_1 \xrightarrow{\gamma} E'
$$
of complexes on $X$ with $E_1$ perfect such that $\beta : E_1 \to E$
restricts to an isomorphism on $U$ and such that
$\alpha = \gamma|_U \circ \beta|_U^{-1}$.
Moreover we can assume $E_1 = E \otimes_{\mathcal{O}_X}^\mathbf{L} I$
for some perfect complex $I$ on $X$.
\end{lemma}

\begin{proof}
Write $X = \Spec(A)$. Write $U = D(f_1) \cup \ldots \cup D(f_r)$. Choose
finite complex of finite projective $A$-modules $M^\bullet$ representing
$E$ (Lemma \href{https://stacks.math.columbia.edu/tag/08EB}{lemma perfect affine (Tag 08EB)}). Choose a complex of $A$-modules
$(M')^\bullet$ representing $E'$ (Lemma \ref{lemma-affine-compare-bounded}).
In this case the complex $H^\bullet = \Hom_A(M^\bullet, (M')^\bullet)$
is a complex of $A$-modules whose associated complex of quasi-coherent
$\mathcal{O}_X$-modules represents $R\SheafHom(E, E')$, see
Cohomology, Lemma \href{https://stacks.math.columbia.edu/tag/08DM}{lemma Rhom strictly perfect (Tag 08DM)}.
Then $\alpha$ determines an element $s$ of $H^0(U, R\SheafHom(E, E'))$, see
Cohomology, Lemma \href{https://stacks.math.columbia.edu/tag/08DK}{lemma section RHom over U (Tag 08DK)}.
There exists an $e$ and a map
$$
\xi : I^\bullet(f_1^e, \ldots, f_r^e) \to \Hom_A(M^\bullet, (M')^\bullet)
$$
corresponding to $s$, see
Proposition \ref{proposition-represent-cohomology-class-on-open}.
Letting $E_1$ be the object corresponding to
complex of quasi-coherent $\mathcal{O}_X$-modules
associated to
$$
\text{Tot}(I^\bullet(f_1^e, \ldots, f_r^e) \otimes_A M^\bullet)
$$
we obtain $E_1 \to E$ using the canonical map
$I^\bullet(f_1^e, \ldots, f_r^e) \to A$ and $E_1 \to E'$
using $\xi$ and
Cohomology, Lemma \href{https://stacks.math.columbia.edu/tag/08DK}{lemma section RHom over U (Tag 08DK)}.
\end{proof}


\begin{lemma}
\label{lemma-lift-perfect-complex-plus-shift}
Let $X$ be an affine scheme. Let $U \subset X$ be a quasi-compact open.
For every perfect object $F$ of $D(\mathcal{O}_U)$
the object $F \oplus F[1]$ is the restriction of
a perfect object of $D(\mathcal{O}_X)$.
\end{lemma}

\begin{proof}
By Lemma \ref{lemma-lift-perfect-complex-plus-locally-free}
we can find a perfect object $E$ of $D(\mathcal{O}_X)$
such that $E|_U = \mathcal{F}[r] \oplus F$ for some finite locally
free $\mathcal{O}_U$-module $\mathcal{F}$.
By Lemma \ref{lemma-lift-map} we can find a morphism of
perfect complexes $\alpha : E_1 \to E$ such that $(E_1)|_U \cong E|_U$
and such that $\alpha|_U$ is the map
$$
\left(
\begin{matrix}
\text{id}_{\mathcal{F}[r]} & 0 \\
0 & 0
\end{matrix}
\right)
:
\mathcal{F}[r] \oplus F \to \mathcal{F}[r] \oplus F
$$
Then the cone on $\alpha$ is a solution.
\end{proof}


\begin{lemma}
\label{lemma-lift-perfect-complex-plus-locally-free}
Let $X$ be an affine scheme. Let $U \subset X$ be a quasi-compact open.
For every perfect object $E$ of $D(\mathcal{O}_U)$ there exists an integer
$r$ and a finite locally free sheaf $\mathcal{F}$ on $U$ such that
$\mathcal{F}[-r] \oplus E$ is the restriction of a perfect object of
$D(\mathcal{O}_X)$.
\end{lemma}

\begin{proof}
Say $X = \Spec(A)$. Recall that a perfect complex is
pseudo-coherent, see
Cohomology, Lemma \ref{cohomology-lemma-perfect}.
By Lemma \ref{lemma-lift-pseudo-coherent} we can find a bounded above complex
$\mathcal{F}^\bullet$ of finite free $A$-modules such that $E$ is
isomorphic to $\mathcal{F}^\bullet|_U$ in $D(\mathcal{O}_U)$.
By Cohomology, Lemma \ref{cohomology-lemma-perfect} and since
$U$ is quasi-compact, we see that $E$ has finite tor dimension, say
$E$ has tor amplitude in $[a, b]$. Pick $r < a$ and set
$$
\mathcal{K} = \Ker(\mathcal{F}^{r} \to \mathcal{F}^{r + 1})
= \Im(\mathcal{F}^{r - 1} \to \mathcal{F}^r).
$$
Since $E$ has tor amplitude in $[a, b]$ we see that
$\mathcal{F} = \mathcal{K}|_U$ is
flat (Cohomology, Lemma \ref{cohomology-lemma-last-one-flat}).
Hence $\mathcal{F}$ is flat and of finite presentation, thus finite
locally free (Properties, Lemma \ref{properties-lemma-finite-locally-free}).
It follows that
$$
\mathcal{F} \to \mathcal{F}^r|_U \to \mathcal{F}^{r + 1}|_U \to \ldots
$$
is a strictly perfect complex on $U$ representing $E$. On the other hand,
the complex $P = (\mathcal{F}^r \to \mathcal{F}^{r + 1} \to \ldots )$
is a perfect complex on $X$. Using
stupid truncations we obtain a distinguished triangle
$$
P|_U \to E \to \mathcal{F}[-r - 1] \to (P|_U)[1]
$$
If the map $E \to \mathcal{F}[-r - 1]$ is zero in $D(\mathcal{O}_U)$,
then $P|_U = \mathcal{F}[-r - 2] \oplus E$, see
Derived Categories, Lemma \href{https://stacks.math.columbia.edu/tag/05QT}{lemma split (Tag 05QT)}.
This will be true for $r \ll 0$ for example by
Lemma \ref{lemma-ext-from-perfect-into-bounded-QCoh}.
\end{proof}


\begin{lemma}
\label{cohomology-lemma-glue}
Let $(X, \mathcal{O}_X)$ be a ringed space. Let $X = U \cup V$ be
the union of two open subspaces of $X$. Suppose given
\begin{enumerate}
\item an object $A$ of $D(\mathcal{O}_U)$,
\item an object $B$ of $D(\mathcal{O}_V)$, and
\item an isomorphism $c : A|_{U \cap V} \to B|_{U \cap V}$.
\end{enumerate}
Then there exists an object $F$ of $D(\mathcal{O}_X)$
and isomorphisms $f : F|_U \to A$, $g : F|_V \to B$ such
that $c = g|_{U \cap V} \circ f^{-1}|_{U \cap V}$.
Moreover, given
\begin{enumerate}
\item an object $E$ of $D(\mathcal{O}_X)$,
\item a morphism $a : A \to E|_U$ of $D(\mathcal{O}_U)$,
\item a morphism $b : B \to E|_V$ of $D(\mathcal{O}_V)$, 
\end{enumerate}
such that
$$
a|_{U \cap V}  = b|_{U \cap V} \circ c.
$$
Then there exists a morphism $F \to E$ in $D(\mathcal{O}_X)$
whose restriction to $U$ is $a \circ f$
and whose restriction to $V$ is $b \circ g$.
\end{lemma}

\begin{proof}
Denote $j_U$, $j_V$, $j_{U \cap V}$ the corresponding open immersions.
Choose a distinguished triangle
$$
F \to Rj_{U, *}A \oplus Rj_{V, *}B \to Rj_{U \cap V, *}(B|_{U \cap V})
\to F[1]
$$
where the map $Rj_{V, *}B \to Rj_{U \cap V, *}(B|_{U \cap V})$ is the
obvious one and where
$Rj_{U, *}A \to Rj_{U \cap V, *}(B|_{U \cap V})$
is the composition of
$Rj_{U, *}A \to Rj_{U \cap V, *}(A|_{U \cap V})$
with $Rj_{U \cap V, *}c$. Restricting to $U$ we obtain
$$
F|_U \to A \oplus (Rj_{V, *}B)|_U \to (Rj_{U \cap V, *}(B|_{U \cap V}))|_U
\to F|_U[1]
$$
Denote $j : U \cap V \to U$. Compatibility of restriction to opens and
cohomology shows that both
$(Rj_{V, *}B)|_U$ and $(Rj_{U \cap V, *}(B|_{U \cap V}))|_U$
are canonically isomorphic to $Rj_*(B|_{U \cap V})$.
Hence the second arrow of the last displayed diagram has
a section, and we conclude that the morphism $F|_U \to A$ is
an isomorphism. Similarly, the morphism $F|_V \to B$ is an
isomorphism. The existence of the morphism $F \to E$ follows
from the Mayer-Vietoris sequence for $\Hom$, see
Lemma \href{https://stacks.math.columbia.edu/tag/08BW}{lemma mayer vietoris hom (Tag 08BW)}.
\end{proof}


\begin{lemma}
\label{cohomology-lemma-pushforward-restriction}
Let $(X, \mathcal{O}_X)$ be a ringed space. Let $j : U \to X$ be an
open subspace. Let $T \subset X$ be a closed subset contained in $U$.
\begin{enumerate}
\item If $E$ is an object of $D(\mathcal{O}_X)$ whose cohomology sheaves
are supported on $T$, then $E \to Rj_*(E|_U)$ is an isomorphism.
\item If $F$ is an object of $D(\mathcal{O}_U)$ whose cohomology sheaves
are supported on $T$, then $j_!F \to Rj_*F$ is an isomorphism.
\end{enumerate}
\end{lemma}

\begin{proof}
Let $V = X \setminus T$ and $W = U \cap V$. Note that $X = U \cup V$ is an
open covering of $X$. Denote $j_W : W \to V$ the open immersion.
Let $E$ be an object of $D(\mathcal{O}_X)$ whose cohomology sheaves are
supported on $T$. By
Lemma \href{https://stacks.math.columbia.edu/tag/08FE}{lemma restrict direct image open (Tag 08FE)} we have
$(Rj_*E|_U)|_V = Rj_{W, *}(E|_W) = 0$ because $E|_W = 0$ by our assumption.
On the other hand, $Rj_*(E|_U)|_U = E|_U$. Thus (1) is clear.
Let $F$ be an object of $D(\mathcal{O}_U)$ whose cohomology sheaves
are supported on $T$. By
Lemma \href{https://stacks.math.columbia.edu/tag/08FE}{lemma restrict direct image open (Tag 08FE)} we have
$(Rj_*F)|_V = Rj_{W, *}(F|_W) = 0$ because $F|_W = 0$ by our assumption.
We also have $(j_!F)|_V = j_{W!}(F|_W) = 0$ (the first equality is immediate
from the definition of extension by zero). Since both
$(Rj_*F)|_U = F$ and $(j_!F)|_U = F$ we see that (2) holds.
\end{proof}


\begin{lemma}
\label{cohomology-lemma-pushforward-perfect}
Let $(X, \mathcal{O}_X)$ be a ringed space. Let $j : U \to X$ be an
open subspace. Let $E$ be a perfect object of $D(\mathcal{O}_U)$
whose cohomology
sheaves are supported on a closed subset $T \subset U$ with $j(T)$
closed in $X$. Then $Rj_*E$ is a perfect object of $D(\mathcal{O}_X)$.
\end{lemma}

\begin{proof}
Being a perfect complex is local on $X$. Thus it suffices to check that
$Rj_*E$ is perfect when restricted to $U$ and $V = X \setminus j(T)$.
We have $Rj_*E|_U = E$ which is perfect. We have
 $Rj_*E|_V = 0$ because $E|_{U \setminus T} = 0$.
\end{proof}


\begin{proposition}
\label{proposition-represent-cohomology-class-on-open}
In Situation \ref{situation-complex}. For every object $E$
of $D_\QCoh(\mathcal{O}_X)$ the map
(\href{https://stacks.math.columbia.edu/tag/08DC}{equation comparison (Tag 08DC)}) is an isomorphism.
\end{proposition}

\begin{proof}
By Lemma \ref{lemma-affine-compare-bounded} we may assume that $E$
is given by a complex of quasi-coherent sheaves $\mathcal{F}^\bullet$.
Let $M^\bullet = \Gamma(X, \mathcal{F}^\bullet)$ be the corresponding
complex of $A$-modules. By
Lemmas \ref{lemma-alternating-cech-complex-complex} and
\ref{lemma-alternating-cech-complex-complex-computes-cohomology}
we have quasi-isomorphisms
$$
\colim_e \Hom^\bullet(I^\bullet(f_1^e, \ldots, f_r^e), M^\bullet)
\longrightarrow
\text{Tot}(\check{\mathcal{C}}_{alt}^\bullet(\mathcal{U}, \mathcal{F}^\bullet))
\longrightarrow
R\Gamma(U, \mathcal{F}^\bullet)
$$
By More on Algebra, Lemma \href{https://stacks.math.columbia.edu/tag/0A66}{lemma RHom out of projective (Tag 0A66)} and
Equation (\href{https://stacks.math.columbia.edu/tag/0A64}{equation h0 RHom (Tag 0A64)})
taking $H^0$ of the complex
$\Hom^\bullet(I^\bullet(f_1^e, \ldots, f_r^e), M^\bullet)$
computes $\Hom$ in $D(A)$. Thus taking $H^0$ on both sides we obtain
$$
\colim_e \Hom_{D(A)}(I^\bullet(f_1^e, \ldots, f_r^e), M^\bullet)
=
H^0(U, E)
$$
Since $\Hom_{D(A)}(I^\bullet(f_1^e, \ldots, f_r^e), M^\bullet) =
\Hom_{D(\mathcal{O}_X)}(I_e, E)$ by
Lemma \ref{lemma-affine-compare-bounded} the lemma follows.
\end{proof}


\begin{lemma}
\label{lemma-represent-cohomology-class-on-closed}
In Situation \ref{situation-complex}. Let $E$ be an object of
$D_\QCoh(\mathcal{O}_X)$.
Assume that $H^i(E)|_U = 0$ for $i = - r + 1, \ldots, 0$.
Then given $s \in H^0(X, E)$ there exists an $e \geq 0$ and
a morphism $K_e \to E$ such that $s$ is in the image of
$H^0(X, K_e) \to H^0(X, E)$.
\end{lemma}

\begin{proof}
Since $U$ is covered by $r$ affine opens we have $H^j(U, \mathcal{F}) = 0$
for $j \geq r$ and any quasi-coherent module
(Cohomology of Schemes, Lemma \href{https://stacks.math.columbia.edu/tag/01XI}{lemma vanishing nr affines (Tag 01XI)}).
By Lemma \ref{lemma-application-nice-K-injective} we see that $H^0(U, E)$
is equal to $H^0(U, \tau_{\geq -r + 1}E)$. There is
a spectral sequence
$$
H^j(U, H^i(\tau_{\geq -r + 1}E)) \Rightarrow H^{i + j}(U, \tau_{\geq -N}E)
$$
see Derived Categories, Lemma \href{https://stacks.math.columbia.edu/tag/015J}{lemma two ss complex functor (Tag 015J)}.
Hence $H^0(U, E) = 0$ by our assumed vanishing of cohomology sheaves of $E$.
We conclude that $s|_U = 0$.
Think of $s$ as a morphism $\mathcal{O}_X \to E$ in $D(\mathcal{O}_X)$.
By Proposition \ref{proposition-represent-cohomology-class-on-open}
the composition $I_e \to \mathcal{O}_X \to E$ is zero for some $e$.
By the distinguished triangle $I_e \to \mathcal{O}_X \to K_e \to I_e[1]$
we obtain a morphism $K_e \to E$ such that $s$ is the composition
$\mathcal{O}_X \to K_e \to E$.
\end{proof}


\begin{lemma}
\label{lemma-alternating-cech-complex}
In Situation \ref{situation-complex}. Let $M$ be an $A$-module and
denote $\mathcal{F}$ the associated $\mathcal{O}_X$-module. Then
there is a canonical isomorphism of complexes
$$
\Psi : \colim_e \Hom_A(I^\bullet(f_1^e, \ldots, f_r^e), M)
\longrightarrow
\check{\mathcal{C}}_{alt}^\bullet(\mathcal{U}, \mathcal{F})
$$
functorial in $M$ where the differentials on the $\Hom$-complex
are the contragredients of the differentials on
$I^\bullet(f_1^e, \ldots, f_r^e)$.
\end{lemma}

\begin{proof}
Recall that the alternating {\v C}ech complex is the subcomplex
of the usual {\v C}ech complex given by alternating cochains, see
Cohomology, Section \href{https://stacks.math.columbia.edu/tag/01FG}{section alternating cech (Tag 01FG)}.
As usual we view a $p$-cochain in
$\check{\mathcal{C}}_{alt}^\bullet(\mathcal{U}, \mathcal{F})$
as an alternating function $s$ on $\{1, \ldots, r\}^{p + 1}$
whose value $s_{i_0\ldots i_p}$ at $(i_0, \ldots, i_p)$ lies in
$M_{f_{i_0}\ldots f_{i_p}} = \mathcal{F}(U_{i_0\ldots i_p})$.
On the other hand, a $p$-cochain $t$ in
$\Hom^\bullet(I^\bullet(f_1^e, \ldots, f_r^e), M)$
is a map $t : \wedge^{p + 1}(A^{\oplus r}) \to M$.
Write $[i] \in A^{\oplus r}$ for the $i$th basis element and
write
$$
[i_0, \ldots, i_p] = [i_0] \wedge \ldots \wedge [i_p]
\in \wedge^{p + 1}(A^{\oplus r})
$$
For $t$ as above we set
$$
\Psi(t)_{i_0 \ldots i_p} =
(-1)^p \frac{t([i_0, \ldots, i_p])}{f_{i_0}^e\ldots f_{i_p}^e}
$$
It is clear that $\Psi(t)$ is an alternating cochain.
The rule above is compatible with the transition maps
of the system as the transition map
$$
I^\bullet(f_1^e, \ldots, f_r^e) \leftarrow
I^\bullet(f_1^{e + 1}, \ldots, f_r^{e + 1}),
$$
of (\ref{equation-system}) sends $[i_0, \ldots, i_p]$
to $f_{i_0}\ldots f_{i_p}[i_0, \ldots, i_p]$.
It is clear from the description of the localizations
$M_{f_{i_0} \ldots f_{i_p}}$ in
Algebra, Lemma \href{https://stacks.math.columbia.edu/tag/00CR}{lemma localization colimit (Tag 00CR)}
that the rule $\Psi$ defines an isomorphism of cochain modules
in degree $p$ in the colimit. To finish the proof we have to show that the map
is compatible with differentials. To see this, for $t$ as above we compute
\begin{align*}
d(\Psi(t))_{i_0 \ldots i_{p + 1}}
& =
\sum\nolimits_{j = 0}^{p + 1} (-1)^j
\Psi(t)_{i_0\ldots \hat i_j \ldots i_{p + 1}} \\
& =
(-1)^p
\sum\nolimits_{j = 0}^{p + 1} (-1)^j t([i_0 \ldots \hat i_j \ldots i_{p + 1}])
(f_{i_0} \ldots \hat f_{i_j} \ldots f_{i_p})^{-e}
\end{align*}
Recall that the differentials on $I^\bullet(f_1^e, \ldots, f_r^e)$
are the negative of the differentials on $K^\bullet(f_1, \ldots, f_r)$.
Thus
\begin{align*}
\Psi(d(t))_{i_0 \ldots i_{p + 1}} & =
(-1)^{p + 1}
d(t)([i_0, \ldots, i_{p + 1}]) (f_{i_0} \ldots f_{i_{p + 1}})^{-e} \\
& =
(-1)^{p + 1}
t(d([i_0, \ldots, i_{p + 1}])) (f_{i_0} \ldots f_{i_{p + 1}})^{-e} \\
& =
(-1)^{p + 1}
t(-\sum\nolimits_{j = 0}^{p + 1}
(-1)^j f_{i_j}^e [i_0, \ldots, \hat i_j, \ldots i_{p + 1}])
(f_{i_0} \ldots f_{i_{p + 1}})^{-e} \\
& =
-(-1)^{p + 1}
\sum\nolimits_{j = 0}^{p + 1}
(-1)^j t([i_0, \ldots, \hat i_j, \ldots i_{p + 1}])
(f_{i_0} \ldots \hat f_{i_j} \ldots f_{i_p})^{-e}
\end{align*}
The two formulas agree concluding the proof.
\end{proof}


\begin{lemma}
\label{lemma-alternating-cech-complex-complex}
In Situation \ref{situation-complex}. Let $M^\bullet$ be a
complex of $A$-modules and
denote $\mathcal{F}^\bullet$ the associated complex of
$\mathcal{O}_X$-modules. Then
there is a canonical isomorphism of complexes
$$
\colim_e \Hom^\bullet(I^\bullet(f_1^e, \ldots, f_r^e), M^\bullet)
\longrightarrow
\text{Tot}(\check{\mathcal{C}}_{alt}^\bullet(\mathcal{U}, \mathcal{F}^\bullet))
$$
functorial in $M^\bullet$.
\end{lemma}

\begin{proof}
Consider the double complex $F^{\bullet, \bullet}$ with terms
$F^{p, q} = \mathcal{C}_{alt}^p(\mathcal{U}, \mathcal{F}^q)$
discussed in
Cohomology, Section \href{https://stacks.math.columbia.edu/tag/01FP}{section cech cohomology of complexes (Tag 01FP)}.
Consider the double complex $G^{\bullet, \bullet}$ with terms
$G^{p, q} = \colim_e \Hom_A(I^{-p}(f_1^e, \ldots, f_r^e), M^q)$ and
differentials given by functoriality (without the intervention
of signs). The maps $\psi^{p, q} : G^{p, q} \to F^{p, q}$
constructed in the proof of Lemma \ref{lemma-alternating-cech-complex}
are isomorphisms and compatible with the differentials $d_1$ (by the lemma)
and $d_2$ (this is clear). However, the differentials $d$ on the
complexes on the left and right hand side of the arrow in the lemma
have different signs. Namely, for $g \in G^{p, q}$ is given by
$$
d(g) =  d_2(g) - (-1)^{p + q} d_1(g) 
$$
(see More on Algebra, Section \href{https://stacks.math.columbia.edu/tag/0A8H}{section hom complexes (Tag 0A8H)})
and the differential for $f \in F^{p, q}$ is given by
$$
d(f) = d_1(f) + (-1)^p d_2(f)
$$
Thus we can fix the signs by multiplying $\psi^{p, q}$ by
$(-1)^{pq + p(p - 1)/2}$.
\end{proof}


\begin{lemma}
\label{lemma-alternating-cech-complex-complex-computes-cohomology}
In Situation \ref{situation-complex}. Let $\mathcal{F}^\bullet$
be a complex of quasi-coherent $\mathcal{O}_X$-modules. Then
there is a canonical isomorphism
$$
\text{Tot}(\check{\mathcal{C}}_{alt}^\bullet(\mathcal{U}, \mathcal{F}^\bullet))
\longrightarrow
R\Gamma(U, \mathcal{F}^\bullet)
$$
in $D(A)$ functorial in $\mathcal{F}^\bullet$.
\end{lemma}

\begin{proof}
Let $\mathcal{B}$ be the set of affine opens of $U$. Since the higher
cohomology groups of a quasi-coherent module on an affine scheme are zero
(Cohomology of Schemes, Lemma
\href{https://stacks.math.columbia.edu/tag/01XB}{lemma quasi coherent affine cohomology zero (Tag 01XB)})
this is a special case of
Cohomology, Lemma \href{https://stacks.math.columbia.edu/tag/08C2}{lemma alternating cech complex complex ss (Tag 08C2)}.
\end{proof}


\begin{theorem}
\label{theorem-approximation}
Let $X$ be a quasi-compact and quasi-separated scheme.
Then approximation by perfect complexes holds on $X$.
\end{theorem}

\begin{proof}
This follows from the induction principle of
Cohomology of Schemes, Lemma \ref{coherent-lemma-induction-principle}
and Lemmas \ref{lemma-induction-step} and \ref{lemma-approximation-affine}.
\end{proof}


\noindent
Let $A$ be a ring and let $f_1, \ldots, f_r$ be a sequence of elements
of $A$. We have defined the Koszul complex
$K_\bullet(f_1, \ldots, f_r)$ in
More on Algebra, Definition \ref{more-algebra-definition-koszul-complex}.
It is a chain complex sitting in degrees $r, \ldots, 0$.
We turn this into a cochain complex $K^\bullet(f_1, \ldots, f_r)$
by setting $K^{-n}(f_1, \ldots, f_r) = K_n(f_1, \ldots, f_r)$
and using the same differentials. In the rest of this section all
the complexes will be cochain complexes.

\medskip\noindent
We define a complex $I^\bullet(f_1, \ldots, f_r)$
such that we have a distinguished triangle
$$
I^\bullet(f_1, \ldots, f_r) \to
A \to
K^\bullet(f_1, \ldots, f_r) \to
I^\bullet(f_1, \ldots, f_r)[1]
$$
in $K(A)$.
In other words, we set
$$
I^i(f_1, \ldots, f_r) =
\left\{
\begin{matrix}
K^{i - 1}(f_1, \ldots, f_r) & \text{if } i \leq 0 \\
0 & \text{else}
\end{matrix}
\right.
$$
and we use the negative of the differential on $K^\bullet(f_1, \ldots, f_r)$.
The maps in the distinguished triangle are the obvious ones. Note that
$I^0(f_1, \ldots, f_r) = A^{\oplus r} \to A$ is given by
multiplication by $f_i$ on the $i$th factor.
Hence $I^\bullet(f_1, \ldots, f_r) \to A$ factors as
$$
I^\bullet(f_1, \ldots, f_r) \to I \to A
$$
where $I = (f_1, \ldots, f_r)$. In fact, there is a short exact sequence
$$
0 \to H^{-1}(K^\bullet(f_1, \ldots, f_s)) \to
H^0(I^\bullet(f_1, \ldots, f_s)) \to I \to 0
$$
and for every $i < 0$ we have
$H^i(I^\bullet(f_1, \ldots, f_r)) = H^{i - 1}(K^\bullet(f_1, \ldots, f_r))$.
Observe that given a second sequence $g_1, \ldots, g_r$ of elements of $A$
there are canonical maps
$$
I^\bullet(f_1g_1, \ldots, f_rg_r) \to I^\bullet(f_1, \ldots, f_r)
\quad\text{and}\quad
K^\bullet(f_1g_1, \ldots, f_rg_r) \to K^\bullet(f_1, \ldots, f_r)
$$
compatible with the maps described above. The first of these maps is
given by multiplication by $g_i$ on the $i$th summand of
$I^0(f_1g_1, \ldots, f_rg_r) = A^{\oplus r}$. In particular, given
$f_1, \ldots, f_r$ we obtain an inverse system of complexes
\begin{equation}
\label{equation-system}
I^\bullet(f_1, \ldots, f_r) \leftarrow
I^\bullet(f_1^2, \ldots, f_r^2) \leftarrow
I^\bullet(f_1^3, \ldots, f_r^3) \leftarrow \ldots
\end{equation}
which will play an important role in that which is to follow.
To easily formulate the following lemmas we fix some notation.


\begin{lemma}
\label{lemma-compute-ext}
Let $f : X \to S$ be a morphism of schemes, $E \in D(\mathcal{O}_X)$
and $\mathcal{G}^\bullet$ a complex of $\mathcal{O}_X$-modules.
Assume
\begin{enumerate}
\item $S$ is Noetherian,
\item $f$ is locally of finite type,
\item $E \in D^-_{\textit{Coh}}(\mathcal{O}_X)$,
\item $\mathcal{G}^\bullet$ is a bounded complex of
coherent $\mathcal{O}_X$-modules flat over $S$ with support proper over $S$.
\end{enumerate}
Then the following two statements are true
\begin{enumerate}
\item[(A)] for every $m \in \mathbf{Z}$ there exists a perfect object $K$
of $D(\mathcal{O}_S)$ and functorial maps
$$
\alpha^i_\mathcal{F} :
\Ext^i_{\mathcal{O}_X}(E,
\mathcal{G}^\bullet \otimes_{\mathcal{O}_X} f^*\mathcal{F})
\longrightarrow
H^i(S, K \otimes^\mathbf{L}_{\mathcal{O}_S} \mathcal{F})
$$
for $\mathcal{F}$ quasi-coherent on $S$ compatible with boundary maps
(see proof) such that $\alpha^i_\mathcal{F}$ is an isomorphism for $i \leq m$
\item[(B)] there exists a pseudo-coherent $L \in D(\mathcal{O}_S)$
and functorial isomorphisms
$$
\Ext^i_{\mathcal{O}_S}(L, \mathcal{F}) \longrightarrow
\Ext^i_{\mathcal{O}_X}(E,
\mathcal{G}^\bullet \otimes_{\mathcal{O}_X} f^*\mathcal{F})
$$
for $\mathcal{F}$ quasi-coherent on $S$ compatible with boundary maps.
\end{enumerate}
\end{lemma}

\begin{proof}
Proof of (A).
Suppose $\mathcal{G}^i$ is nonzero only for $i \in [a, b]$.
We may replace $X$ by a quasi-compact open neighbourhood of
the union of the supports of $\mathcal{G}^i$.
Hence we may assume $X$ is Noetherian.
In this case $X$ and $f$ are quasi-compact and quasi-separated.
Choose an approximation $P \to E$ by a perfect complex $P$ of
$(X, E, -m - 1 + a)$
(possible by Theorem \ref{theorem-approximation}).
Then the induced map
$$
\Ext^i_{\mathcal{O}_X}(E,
\mathcal{G}^\bullet \otimes_{\mathcal{O}_X} f^*\mathcal{F})
\longrightarrow
\Ext^i_{\mathcal{O}_X}(P,
\mathcal{G}^\bullet \otimes_{\mathcal{O}_X} f^*\mathcal{F})
$$
is an isomorphism for $i \leq m$. Namely, the kernel, resp.\ cokernel of this
map is a quotient, resp.\ submodule of
$$
\Ext^i_{\mathcal{O}_X}(C,
\mathcal{G}^\bullet \otimes_{\mathcal{O}_X} f^*\mathcal{F})
\quad\text{resp.}\quad
\Ext^{i + 1}_{\mathcal{O}_X}(C,
\mathcal{G}^\bullet \otimes_{\mathcal{O}_X} f^*\mathcal{F})
$$
where $C$ is the cone of $P \to E$. Since $C$ has vanishing cohomology
sheaves in degrees $\geq -m - 1 + a$ these $\Ext$-groups are zero
for $i \leq m + 1$ by
Derived Categories, Lemma \href{https://stacks.math.columbia.edu/tag/06XS}{lemma negative exts (Tag 06XS)}.
This reduces us to the case that
$E$ is a perfect complex which is Lemma \ref{lemma-compute-ext-perfect}.
The statement on boundaries is explained in the proof of
Lemma \ref{lemma-compute-ext-perfect}.

\medskip\noindent
Proof of (B). As in the proof of (A) we may assume $X$ is Noetherian.
Observe that $E$ is pseudo-coherent by
Lemma \href{https://stacks.math.columbia.edu/tag/08E8}{lemma identify pseudo coherent noetherian (Tag 08E8)}.
By Lemma \ref{lemma-pseudo-coherent-hocolim} we can write
$E = \text{hocolim} E_n$ with $E_n$ perfect and $E_n \to E$ inducing
an isomorphism on truncations $\tau_{\geq -n}$. Let $E_n^\vee$
be the dual perfect complex
(Cohomology, Lemma \href{https://stacks.math.columbia.edu/tag/08DQ}{lemma dual perfect complex (Tag 08DQ)}).
We obtain an inverse system $\ldots \to E_3^\vee \to E_2^\vee \to E_1^\vee$
of perfect objects. This in turn gives rise to an inverse system
$$
\ldots \to K_3 \to K_2 \to K_1\quad\text{with}\quad
K_n = Rf_*(E_n^\vee \otimes_{\mathcal{O}_X}^\mathbf{L} \mathcal{G}^\bullet)
$$
perfect on $S$, see Lemma \ref{lemma-tensor-perfect}.
By Lemma \ref{lemma-compute-ext-perfect} and its proof and
by the arguments in the previous paragraph (with $P = E_n$)
for any quasi-coherent $\mathcal{F}$ on $S$ we have
functorial canonical maps
$$
\xymatrix{
& \Ext^i_{\mathcal{O}_X}(E,
\mathcal{G}^\bullet \otimes_{\mathcal{O}_X} f^*\mathcal{F})
\ar[ld] \ar[rd] \\
H^i(S, K_{n + 1} \otimes_{\mathcal{O}_S}^\mathbf{L} \mathcal{F})
\ar[rr] & &
H^i(S, K_n \otimes_{\mathcal{O}_S}^\mathbf{L} \mathcal{F})
}
$$
which are isomorphisms for $i \leq n + a$.
Let $L_n = K_n^\vee$ be the dual perfect complex.
Then we see that $L_1 \to L_2 \to L_3 \to \ldots$
is a system of perfect objects in $D(\mathcal{O}_S)$
such that for any quasi-coherent $\mathcal{F}$ on $S$
the maps
$$
\Ext^i_{\mathcal{O}_S}(L_{n + 1}, \mathcal{F})
\longrightarrow
\Ext^i_{\mathcal{O}_S}(L_n, \mathcal{F})
$$
are isomorphisms for $i \leq n + a - 1$. This implies that
$L_n \to L_{n + 1}$ induces an isomorphism on truncations
$\tau_{\geq -n - a + 2}$ (hint: take cone of $L_n \to L_{n + 1}$
and look at its last nonvanishing cohomology sheaf).
Thus $L = \text{hocolim} L_n$ is pseudo-coherent, see
Lemma \ref{lemma-pseudo-coherent-hocolim}. The mapping property
of homotopy colimits gives that
$\Ext^i_{\mathcal{O}_S}(L, \mathcal{F}) =
\Ext^i_{\mathcal{O}_S}(L_n, \mathcal{F})$
for $i \leq n + a - 3$ which finishes the proof.
\end{proof}
\end{document}
